\documentclass[11pt]{report}

\usepackage[utf8]{inputenc}
\usepackage[T1]{fontenc}
\usepackage{amsmath,amssymb,amsthm}
\usepackage{booktabs}
\usepackage{longtable}
\usepackage{array}
\usepackage[margin=1in]{geometry}
\usepackage{hyperref}
\usepackage{xcolor}
\usepackage{caption}
\usepackage{enumitem}
\usepackage{graphicx}
\usepackage{float}

\hypersetup{colorlinks=true, linkcolor=blue, citecolor=blue, urlcolor=blue}

\newcommand{\pkg}[1]{\texttt{#1}}
\theoremstyle{plain}
\newtheorem{theorem}{Theorem}[chapter]
\newtheorem{lemma}[theorem]{Lemma}
\newtheorem{proposition}[theorem]{Proposition}
\newtheorem{corollary}[theorem]{Corollary}
\theoremstyle{definition}
\newtheorem{definition}[theorem]{Definition}
\newtheorem{assumption}[theorem]{Assumption}
\theoremstyle{remark}
\newtheorem{remark}[theorem]{Remark}
\renewcommand{\qedsymbol}{$\blacksquare$}

\title{\textbf{Research and Teaching Extensions to the PSAT 2.1.11\\
Power System Analysis Toolbox}\\[0.3em]
\large A Technical Report on Architecture, Mathematical Formulations,\\
Verification Methodology, and Test Results}
\author{Prof. Robert A. Sowah\\
\small Control and Optimization in Energy Networks\\
\small Power Systems Laboratory, PSAT-Based Toolbox Extensions Project}
\date{19 October 2026}

\begin{document}

\maketitle

\begin{abstract}
\noindent
PSAT 2.1.11 (Power System Analysis Toolbox, F.\ Milano) is a mature
MATLAB-based research tool for power-flow, continuation power-flow,
optimal power-flow, and time-domain/small-signal dynamic analysis. As
shipped, it has no documented data schema, no structured results API, no
distributed or uncertainty-aware optimization capability, no storage or
grid-forming-inverter device model, no cyberattack or market-economics
tooling, and no machine-learning or reinforcement-learning interface.
This report documents a coordinated set of extensions -- 25 Python
modules, 31 MATLAB functions, and two genuine new PSAT device classes
registered into PSAT's own class system (\pkg{@ESSclass},
\pkg{@GFMclass}) -- built to close these gaps for graduate teaching and
for modern power-systems research, without modifying any of PSAT's
existing, working behaviour. Every capability reported here was verified
against an independent method, a known closed-form result, or a live
PSAT/MATLAB session before being accepted; this report documents not
only the final, working mathematics but, where instructive, the genuine
defects found and corrected during development, since several of them
recur across more than one module and are worth recording as general
lessons. Beyond stating each extension's governing equations, this
report develops their theoretical basis formally: a dedicated
Mathematical Preliminaries chapter establishes the convex-duality,
fixed-point, graph-spectral, probabilistic, and optimal-control results
every later chapter draws on, and each subsequent chapter states its
own load-bearing claims as numbered Theorems, Lemmas, Propositions, and
Corollaries with proofs -- 36 such results in total, including a formal
Bertsimas--Sim robust-counterpart derivation, an exact two-directional
characterization of undetectable false-data-injection attacks, a
Sherman--Morrison derivation of the line-outage distribution factor, a
Riccati-recursion proof that receding-horizon MPC converges to the
infinite-horizon LQR gain, and a Routh--Hurwitz derivation of the
multi-machine small-signal $90^\circ$ instability threshold --
referencing standard graduate texts (Boyd \& Vandenberghe, Bertsekas,
Horn \& Johnson, Abur \& Exp\'osito, Kundur, and others) throughout,
rather than asserting results without derivation. The report is
organised by capability area: mathematical preliminaries, core
workflow/schema extensions, deep-integration device classes, distribution
and fault analysis, optimization and dispatch under uncertainty, security
and resilience, control, learning-based methods, uncertainty
quantification, and MATPOWER interoperability. A full nomenclature table,
rigorously derived mathematical formulations, and consolidated
test-case results accompany each area.
\end{abstract}

\tableofcontents

%=============================================================
\chapter{Introduction}
%=============================================================

\section{Motivation}

PSAT 2.1.11 is widely used for teaching power-flow, continuation
power-flow (voltage-stability), optimal power-flow, and
time-domain/small-signal dynamic analysis, and it ships with roughly
forty ready-made test systems. It is, however, a \emph{research} tool
first: its device classes (\pkg{@SYclass}, \pkg{@EXclass},
\pkg{@TGclass}, \pkg{@PQclass}, \ldots) have undocumented column
schemas; its results are scattered across global variables with no
structured API; its optimization capability (\texttt{opf}) is a single
deterministic, centralized solve; and it has no native model at all for
battery storage, grid-forming inverters, gas-electric coupling, market
economics, cyberattack surfaces, or learning-based surrogates --
capabilities now standard in the modern power-systems research
literature and in this course's own 16-week syllabus.

This project closes each of those gaps with a self-contained extension,
built and verified in either Python (as a reusable library depending
only on \pkg{numpy}/\pkg{scipy}/\pkg{cvxpy}/\pkg{networkx}/\pkg{torch}/
\pkg{gymnasium}) or MATLAB (as a function operating on PSAT's own loaded
data, or as a PSAT-native device class), or both.

\section{Design principles}

Four principles were applied uniformly across every extension in this
report:

\begin{enumerate}
\item \textbf{Additive, non-destructive integration.} No existing PSAT
  behaviour was changed. The two new device classes
  (\S\ref{ch:deviceclasses}) were wired into PSAT's core dispatch files
  by \emph{adding} lines that reference the new device type; every
  original file was backed up before editing.
\item \textbf{Test before embed.} Every numeric claim in this report was
  produced by code that was actually executed, not assumed correct by
  inspection. Wherever an independent cross-check existed (a brute-force
  enumeration, a different solver, a closed-form textbook result, a
  second programming language), it was used.
\item \textbf{Honest reporting of negative and null results.} Several
  extensions report a finding that did \emph{not} go the way a first
  intuition expected (\S\ref{ch:learning}, \S\ref{ch:resilience}); these
  are reported as found, not adjusted to look better.
\item \textbf{Cross-language parity.} Every Python module with a natural
  PSAT-side analogue has a MATLAB port, tested live against a running
  MATLAB R2026a/PSAT 2.1.11 session (not merely written to look
  plausible). Table~\ref{tab:parity} in \S\ref{ch:verification} gives the
  complete parity map.
\end{enumerate}

\section{Report structure}

Chapter~\ref{ch:nomenclature} is a complete nomenclature table.
Chapter~\ref{ch:prelim} collects, with full proofs or precise
citations, every piece of convex-optimization, fixed-point, graph-linear-
algebra, probability, and optimal-control theory invoked more than once
later in the report, so that each later chapter can cite a numbered
result rather than re-derive it. Chapter~\ref{ch:architecture} describes
the data schema and I/O layer every other extension builds on.
Chapters~\ref{ch:deviceclasses}--\ref{ch:interop} cover each capability
area in turn: mathematical formulation (stated, where the result is
load-bearing, as a numbered Theorem/Lemma/Proposition with a proof, not
only an equation), algorithm, and verified test-case results. Every
Theorem, Lemma, Proposition, and Corollary in the report shares one
numbering sequence per chapter (e.g.\ Theorem~\ref{thm:bertsimas-sim} is
the one and only ``Theorem 7.6''), so a forward or backward
cross-reference always resolves unambiguously.
Chapter~\ref{ch:verification} gives the cross-language parity map and a
consolidated log of every genuine defect found during development -- a
resource in its own right, since several of these defects
(sign-convention bugs, slack-bus formulation errors, nested-function
closures) recur in predictable, nameable ways. Chapter~\ref{ch:limitations}
states every known limitation plainly. Chapter~\ref{ch:conclusion}
concludes.

%=============================================================
\chapter{Nomenclature}
\label{ch:nomenclature}
%=============================================================

Symbols are grouped by the subject area in which they are used; several
letters (e.g.\ $H$, $K$) are reused across areas with different meanings,
as is standard in the power-systems literature -- the local meaning is
always restated at first use within each section.

\section{Network and power-flow quantities}
\begin{longtable}{@{}p{3cm}p{11.5cm}@{}}
\toprule
Symbol & Meaning \\
\midrule
\endhead
$n$, $m$ & Number of buses; number of branches/lines. \\
$\theta_i$ & Voltage phase angle at bus $i$ (rad). \\
$V_i$ & Voltage magnitude at bus $i$ (pu). \\
$B$, $B'$ & Full and reduced (non-slack) nodal susceptance (Laplacian-like) matrix. \\
$Y_{\text{bus}}$, $Z_{\text{bus}}$ & Bus admittance matrix and its inverse, the bus impedance matrix. \\
$P_{g,i}$, $P_{d,i}$ & Real power generation and demand at bus $i$. \\
$\text{PTDF}$ & Power Transfer Distribution Factor matrix, $\partial(\text{flow})/\partial(\text{injection})$. \\
$\text{LODF}_{l,k}$ & Line Outage Distribution Factor: fraction of outaged line $k$'s pre-contingency flow redistributed onto monitored line $l$. \\
$\text{rate}A$ & Thermal flow rating of a line (MW or pu). \\
\bottomrule
\end{longtable}

\section{Optimization quantities}
\begin{longtable}{@{}p{3cm}p{11.5cm}@{}}
\toprule
Symbol & Meaning \\
\midrule
\endhead
$a_i, b_i$ & Quadratic and linear generator cost coefficients, $\text{cost}_i(P)=a_iP^2+b_iP$. \\
$\lambda$ & Lagrange multiplier / shadow price of a power-balance equality constraint (the system marginal cost, or locational marginal price, LMP). \\
$\Gamma$ & Bertsimas--Sim ``budget of uncertainty'' (number of uncertain parameters simultaneously allowed to reach their worst case). \\
$\hat{d}_i$ & Known deviation magnitude of uncertain parameter $i$; true value $d_i=d_i^{\text{nom}}+\xi_i\hat d_i$, $\xi_i\in[-1,1]$. \\
$\rho$ & ADMM penalty parameter. \\
$z$, $u$ & ADMM global consensus variable and scaled dual variable. \\
$W$ & Metropolis--Hastings consensus weight matrix (symmetric, doubly stochastic). \\
$\text{VSS}$ & Value of the Stochastic Solution (expected-value solution cost minus two-stage recourse cost). \\
\bottomrule
\end{longtable}

\section{Dynamics and control quantities}
\begin{longtable}{@{}p{3cm}p{11.5cm}@{}}
\toprule
Symbol & Meaning \\
\midrule
\endhead
$H_i$ & Inertia constant of machine/virtual machine $i$ (s). \\
$D_i$ & Damping coefficient of machine/area $i$ (pu/pu). \\
$\omega$ & Per-unit rotor (or virtual-rotor) speed, nominal 1. \\
$K_{ij}(\theta^\star)$ & Synchronizing-power (small-signal stiffness) coefficient, $B_{ij}\cos(\theta_i^\star-\theta_j^\star)$. \\
$A,B,C,D$ (state space) & Continuous-time linear state-space matrices, $\dot x = Ax+Bu$, $y=Cx+Du$. \\
$A_d, B_d$ & Zero-order-hold discretized equivalents of $A, B$. \\
$Q,R_w$ & State and input quadratic-cost weighting matrices for LQR/MPC. \\
$K_{\text{lqr}}$ & Infinite-horizon LQR gain, from the discrete algebraic Riccati equation (DARE). \\
$\text{ACE}_k$ & Area control error of area $k$: tie-line flow deviation plus frequency-bias term. \\
\bottomrule
\end{longtable}

\section{State estimation and security quantities}
\begin{longtable}{@{}p{3cm}p{11.5cm}@{}}
\toprule
Symbol & Meaning \\
\midrule
\endhead
$H$ (measurement) & Measurement Jacobian, $m\times n$ ($m$ measurements, $n$ state variables). \\
$z$ (measurement) & Measurement vector. \\
$R$ & Measurement noise covariance (diagonal). \\
$\hat x$ & Weighted-least-squares state estimate. \\
$J$ & Chi-square bad-data detection statistic, $J=r^\top R^{-1}r$. \\
$a$ & Measurement attack vector; a \emph{stealthy} attack satisfies $a=Hc$ for some state-shift vector $c$. \\
\bottomrule
\end{longtable}

\section{Gas-network and thermal-rating quantities}
\begin{longtable}{@{}p{3cm}p{11.5cm}@{}}
\toprule
Symbol & Meaning \\
\midrule
\endhead
$p_i$ & Nodal gas pressure (pressure plays the role voltage plays in power flow). \\
$K_{ij}$ (Weymouth) & Weymouth constant of pipe $(i,j)$. \\
$q_{ij}$ & Steady-state gas flow on pipe $(i,j)$. \\
$I_{\max}$ & Dynamic conductor ampacity (A), from the IEEE 738 heat balance. \\
$q_c, q_r, q_s$ & Convective, radiative, and solar heat flux terms (W/m). \\
$R(T_c)$ & Conductor AC resistance at conductor temperature $T_c$. \\
\bottomrule
\end{longtable}

\section{Learning-based-method quantities}
\begin{longtable}{@{}p{3cm}p{11.5cm}@{}}
\toprule
Symbol & Meaning \\
\midrule
\endhead
$x_v$ & Node feature vector (GNN). \\
$e_{uv}$ & Edge feature vector (GNN). \\
$\text{msg}_{uv}$ & Message computed on directed edge $u\to v$. \\
$\text{agg}_v$ & Sum-aggregated incoming messages at node $v$. \\
$s_t, a_t, r_t$ & RL state, action, and reward at time $t$. \\
$\gamma_{\text{RL}}$ & Reinforcement-learning discount factor. \\
\bottomrule
\end{longtable}

%=============================================================
\chapter{Mathematical Preliminaries}
\label{ch:prelim}
%=============================================================

Every later chapter invokes one or more of five bodies of theory:
convex optimization and duality, fixed-point/contraction arguments,
structured linear algebra on graphs, probability theory for Gaussian
and quadratic-form random variables, and discrete-time optimal control.
This chapter states each result once, in the generality needed later,
with a proof or a precise citation to a standard graduate reference, so
that every subsequent chapter can invoke it by number rather than
re-deriving it. Readers already fluent in this material may skip to
Chapter~\ref{ch:architecture}; readers who want full proofs beyond the
sketches given here are pointed to the cited source at each result.

\section{Convex optimization, Lagrangian duality, and KKT conditions}
\label{sec:convex-prelim}

\begin{definition}[Convex set and convex function]
A set $C\subseteq\mathbb{R}^n$ is \emph{convex} if
$\lambda x+(1-\lambda)y\in C$ for all $x,y\in C$, $\lambda\in[0,1]$. A
function $f:C\to\mathbb{R}$ is convex if its epigraph is convex,
equivalently $f(\lambda x+(1-\lambda)y)\le \lambda f(x)+(1-\lambda)f(y)$
for all $x,y\in C$, $\lambda\in[0,1]$; $f$ is \emph{strictly convex} if
the inequality is strict whenever $x\ne y$, $\lambda\in(0,1)$.
\end{definition}

\begin{definition}[Primal problem and Lagrangian]
\label{def:primal}
Consider
\begin{equation}
\min_{x\in\mathbb{R}^n}\; f_0(x) \quad\text{s.t.}\quad g_i(x)\le0 \;(i=1,\ldots,m), \quad h_j(x)=0 \;(j=1,\ldots,p),
\label{eq:primal}
\end{equation}
with $f_0,g_i$ convex and $h_j$ affine. The Lagrangian is
$L(x,\mu,\nu)=f_0(x)+\sum_i\mu_ig_i(x)+\sum_j\nu_jh_j(x)$, $\mu\ge0$, and
the dual function is $d(\mu,\nu)=\inf_x L(x,\mu,\nu)$.
\end{definition}

\begin{assumption}[Slater's condition]
\label{as:slater}
There exists a strictly feasible point $\tilde x$ with $g_i(\tilde
x)<0$ for every $i$ and $h_j(\tilde x)=0$ for every $j$.
\end{assumption}

\begin{theorem}[Strong duality and KKT necessity]
\label{thm:strong-duality}
Under Assumption~\ref{as:slater}, if $x^\star$ solves the primal
problem above then there exist $\mu^\star\ge0,\nu^\star$ (dual
optimal, attaining the primal optimal value with zero duality gap) such
that the following Karush--Kuhn--Tucker (KKT) conditions hold:
\begin{gather}
\nabla f_0(x^\star)+\textstyle\sum_i\mu_i^\star\nabla g_i(x^\star)+\sum_j\nu_j^\star\nabla h_j(x^\star) = 0 \quad\text{(stationarity)}\\
g_i(x^\star)\le0,\; h_j(x^\star)=0 \quad\text{(primal feasibility)}\\
\mu_i^\star\ge0 \quad\text{(dual feasibility)}\\
\mu_i^\star g_i(x^\star)=0 \;\; \forall i \quad\text{(complementary slackness)}
\end{gather}
Conversely, for a convex problem, any $(x,\mu,\nu)$ satisfying the KKT
conditions is primal-dual optimal (sufficiency).
\end{theorem}
\begin{proof}[Proof sketch]
Slater's condition rules out a duality gap for convex problems (Boyd \&
Vandenberghe~\cite{boydconvex}, Ch.~5): strong duality
$f_0(x^\star)=d(\mu^\star,\nu^\star)$ holds, and since
$d(\mu^\star,\nu^\star)=\inf_x L(x,\mu^\star,\nu^\star)\le
L(x^\star,\mu^\star,\nu^\star)=f_0(x^\star)+\sum_i\mu_i^\star
g_i(x^\star)\le f_0(x^\star)$ (the last step using $\mu^\star\ge0$,
$g_i(x^\star)\le0$), equality throughout forces both the complementary
slackness condition above and that $x^\star$ minimizes
$L(\cdot,\mu^\star,\nu^\star)$, which for differentiable $f_0,g_i,h_j$
gives the stationarity condition above. Sufficiency for convex problems
follows because the KKT stationarity condition makes $x$ a global
minimizer of the convex function $L(\cdot,\mu,\nu)$, and complementary
slackness then equates $L(x,\mu,\nu)$ with $f_0(x)$.
\end{proof}

\begin{remark}
Theorem~\ref{thm:strong-duality} is the formal basis for every dual
variable interpreted as a price or shadow price in this report (the LMP
of \S\ref{sec:merit-order-theory}, the reserve price of
\S\ref{sec:chance-theory}, the ADMM/dual-decomposition prices of
\S\ref{sec:decomp-theory}): complementary slackness is exactly why a
non-binding constraint has zero price and a binding constraint's price
equals the constraint's marginal effect on the optimal objective (the
envelope-theorem reading of the stationarity condition).
\end{remark}

\section{Fixed-point and contraction-mapping results}
\label{sec:fixedpoint-prelim}

\begin{definition}[Contraction]
A map $T:X\to X$ on a complete metric space $(X,d)$ is a
\emph{contraction} with modulus $0\le\gamma<1$ if
$d(T(x),T(y))\le\gamma\, d(x,y)$ for all $x,y\in X$.
\end{definition}

\begin{theorem}[Banach fixed-point theorem]
\label{thm:banach}
If $T:X\to X$ is a contraction on a complete metric space, then $T$ has
a unique fixed point $x^\star=T(x^\star)$, and for any $x_0\in X$ the
iterates $x_{k+1}=T(x_k)$ converge to $x^\star$ geometrically,
$d(x_k,x^\star)\le\gamma^k d(x_0,x^\star)$.
\end{theorem}
\begin{proof}
Standard; see Bertsekas~\cite{bertsekasDP}, Appendix~B, or any real
analysis text. Existence/uniqueness follow from showing $(x_k)$ is
Cauchy (via the triangle inequality and $\sum\gamma^k<\infty$) and that
the limit is fixed by continuity of $T$; the rate bound follows by
induction on $d(x_{k+1},x^\star)=d(T(x_k),T(x^\star))\le\gamma\,
d(x_k,x^\star)$.
\end{proof}

\begin{remark}
Theorem~\ref{thm:banach} underlies two independent results used later:
the Bellman-operator contraction behind tabular Q-learning/value
iteration (\S\ref{sec:rl-theory}), and, via the spectral-radius argument
of \S\ref{sec:graph-prelim}, the geometric convergence of Metropolis--
Hastings consensus averaging (\S\ref{sec:decomp-theory}).
\end{remark}

\section{Structured linear algebra on graphs}
\label{sec:graph-prelim}

\begin{definition}[Weighted graph Laplacian]
For a connected graph with $n$ nodes, edge admittances $y_{ij}>0$, the
(weighted) Laplacian $B\in\mathbb{R}^{n\times n}$ has
$B_{ii}=\sum_{j\sim i}y_{ij}$, $B_{ij}=-y_{ij}$ for $i\sim j$, $0$
otherwise.
\end{definition}

\begin{theorem}[Laplacian rank and the reduced system]
\label{thm:laplacian}
$B$ is symmetric positive semi-definite with exactly one zero
eigenvalue (eigenvector $\mathbf{1}$) if and only if the graph is
connected. Consequently, deleting any one row/column (the reference
bus) from $B$ gives a principal submatrix $B'$ that is symmetric
positive definite, hence invertible.
\end{theorem}
\begin{proof}
For any $x\in\mathbb{R}^n$, $x^\top Bx=\sum_{i\sim j}y_{ij}(x_i-x_j)^2\ge0$,
so $B\succeq0$; equality holds iff $x_i=x_j$ for every edge, i.e.\ iff
$x$ is constant on each connected component, giving a null space of
dimension equal to the number of components (one, under connectivity) --
this is the standard Kirchhoff/Matrix-Tree argument (see Horn \&
Johnson~\cite{hornjohnson}, \S13.4, or Godsil \& Royle, \emph{Algebraic
Graph Theory}). A symmetric PSD matrix with a single zero eigenvalue has
every principal submatrix obtained by deleting one row/column strictly
positive definite, because the deleted-coordinate quadratic form
$x'^\top B'x'$ (extending $x'$ by a zero in the deleted coordinate) can
only vanish for $x'=0$, since the unique null direction $\mathbf 1$ has a
nonzero entry at the deleted coordinate.
\end{proof}

\begin{remark}
Theorem~\ref{thm:laplacian} is why every DC-network formulation in this
report (Chapters~\ref{ch:optimization}, \ref{ch:resilience}) fixes one
reference-bus angle and solves only the reduced system $B'\theta'=P'$: it
is the formal reason $B'$ is guaranteed nonsingular whenever the network
is connected, matching the main course's own Theorem~1 treatment of the
identical fact for AC power flow's own $Y_{\text{bus}}$-derived
Jacobian structure.
\end{remark}

\begin{theorem}[Sherman--Morrison rank-one update]
\label{thm:sherman-morrison}
For invertible $A\in\mathbb{R}^{n\times n}$ and $u,v\in\mathbb{R}^n$
with $1+v^\top A^{-1}u\ne0$,
\begin{equation}
(A+uv^\top)^{-1} = A^{-1} - \frac{A^{-1}uv^\top A^{-1}}{1+v^\top A^{-1}u}.
\label{eq:sherman-morrison}
\end{equation}
\end{theorem}
\begin{proof}
Multiply the right-hand side by $(A+uv^\top)$ and simplify using
$v^\top A^{-1}u$ as a scalar; the cross terms telescope to the identity.
See Horn \& Johnson~\cite{hornjohnson}, \S0.7.4, or Golub \& Van
Loan~\cite{golubvanloan}, \S2.1.3.
\end{proof}

\begin{theorem}[Doubly stochastic consensus convergence]
\label{thm:doubly-stochastic}
Let $W\in\mathbb{R}^{n\times n}$ be symmetric, doubly stochastic
($W\mathbf 1=\mathbf 1$, $W=W^\top$, $W\ge0$ entrywise), and
irreducible (the underlying communication graph is connected). Then the
eigenvalues of $W$ satisfy $1=\mu_1>\mu_2\ge\cdots\ge\mu_n\ge-1$, and
for any $x_0\in\mathbb{R}^n$,
\begin{equation}
W^k x_0 \;\longrightarrow\; \bar x\,\mathbf 1, \qquad \bar x=\tfrac1n\textstyle\sum_i x_{0,i},
\label{eq:consensus-limit}
\end{equation}
at the geometric rate $\|W^kx_0-\bar x\mathbf 1\|_2\le \left|\mu_2\right|^k\|x_0-\bar
x\mathbf1\|_2$.
\end{theorem}
\begin{proof}
By the spectral theorem, symmetric $W$ has an orthonormal eigenbasis;
$\mathbf1$ is an eigenvector with eigenvalue $1$ (row sums are $1$). Since
$W$ is doubly stochastic its spectral radius is $1$ (a stochastic
matrix's eigenvalues lie in the closed unit disc by the
Gershgorin circle theorem applied to its rows), and irreducibility plus
the Perron--Frobenius theorem (see Horn \& Johnson~\cite{hornjohnson},
Ch.~8, or Xiao \& Boyd~\cite{xiaoboydconsensus}) make $1$ a simple
eigenvalue with every other eigenvalue satisfying $|\lambda|<1$
whenever $W$ is also aperiodic (guaranteed here by
positive diagonal entries, $W_{ii}>0$, from the Metropolis--Hastings
construction). Decomposing $x_0=\bar x\mathbf1+e_0$ with $e_0\perp
\mathbf1$, $W^kx_0=\bar x\mathbf1+W^ke_0$ and $\|W^ke_0\|_2\le
\left|\mu_2\right|^k\|e_0\|_2\to0$.
\end{proof}

\begin{remark}
Theorem~\ref{thm:doubly-stochastic} is the exact reason this project's
Metropolis--Hastings weight construction (\S\ref{sec:decomp-theory}) is
required to be symmetric and doubly stochastic, not merely row-
stochastic: a row-stochastic-only matrix generally has a LEFT Perron
eigenvector $\pi\ne\mathbf1/n$, so $W^kx_0\to(\pi^\top x_0)\mathbf1$, the
$\pi$-weighted average, not the arithmetic mean -- precisely the bug this
project's own development documents (\S\ref{sec:decomp-theory},
Table~\ref{tab:bugs}) for both \pkg{decentralized\_\hspace{0pt}opt.py}
and \pkg{@GFMclass}/\pkg{@ESSclass}'s consensus-dispatch counterpart.
\end{remark}

\section{Probability theory for Gaussian and quadratic-form random variables}
\label{sec:prob-prelim}

\begin{theorem}[Affine transformation of a Gaussian vector]
\label{thm:affine-gaussian}
If $X\sim\mathcal N(\mu,\Sigma)$ and $Y=AX+b$ for fixed $A\in
\mathbb{R}^{m\times n}$, $b\in\mathbb{R}^m$, then
$Y\sim\mathcal{N}(A\mu+b,\,A\Sigma A^\top)$.
\end{theorem}
\begin{proof}
The characteristic function of $X$ is
$\varphi_X(t)=\exp(it^\top\mu-\tfrac12t^\top\Sigma t)$; that of $Y$ is
$\varphi_Y(t)=\mathbb E[e^{it^\top(AX+b)}]=e^{it^\top
b}\varphi_X(A^\top t)=\exp\left(it^\top(A\mu+b)-\tfrac12t^\top
(A\Sigma A^\top)t\right)$, which is exactly the characteristic function of
$\mathcal N(A\mu+b,A\Sigma A^\top)$; a distribution is determined by its
characteristic function (see Casella \& Berger~\cite{casellaberger},
Thm.~4.6.4, or any graduate probability text).
\end{proof}

\begin{theorem}[Distribution of a quadratic form in Gaussian residuals]
\label{thm:chisquare}
If $r\sim\mathcal N(0,\Sigma_r)$ with $\Sigma_r\succ0$, then
$J=r^\top\Sigma_r^{-1}r\sim\chi^2_k$ for some $k$ depending on the rank
structure of $\Sigma_r$; in particular, if $r=z-H\hat x$ is the
weighted-least-squares residual of an $m$-measurement, $n$-state linear
model with $\hat x=(H^\top R^{-1}H)^{-1}H^\top R^{-1}z$ and
$\Sigma_r=R-H(H^\top R^{-1}H)^{-1}H^\top$, then $J\sim\chi^2_{m-n}$.
\end{theorem}
\begin{proof}
Write $\Sigma_r=LL^\top$ (Cholesky); then $L^{-1}r\sim\mathcal
N(0,I_k)$ (Theorem~\ref{thm:affine-gaussian}) and $J=\|L^{-1}r\|_2^2$ is
a sum of $k$ independent squared standard normals, the definition of
$\chi^2_k$. For the WLS residual, $r=(I-H(H^\top
R^{-1}H)^{-1}H^\top R^{-1})z$ is a linear projection of the Gaussian
measurement noise onto an $(m-n)$-dimensional subspace orthogonal (in the
$R^{-1}$ inner product) to $\text{range}(H)$, giving $k=m-n$ degrees of
freedom; see Abur \& Exposito~\cite{aburexposito}, Ch.~5, for the full
power-systems-state-estimation derivation.
\end{proof}

\begin{remark}
Theorem~\ref{thm:chisquare} is the formal basis for
\S\ref{sec:fdia-theory}'s chi-square bad-data test and, via
Theorem~\ref{thm:affine-gaussian} directly, for
\S\ref{sec:ppf-theory}'s exact analytical probabilistic power flow
($\Sigma_{\text{flow}}=\text{PTDF}\,\Sigma_{\text{inj}}\,
\text{PTDF}^\top$ is nothing but Theorem~\ref{thm:affine-gaussian}
applied to the linear map $\text{flow}=\text{PTDF}\cdot
\text{injections}$).
\end{remark}

\section{Discrete-time optimal control}
\label{sec:control-prelim}

\begin{theorem}[Bellman's principle of optimality and the discrete Riccati recursion]
\label{thm:riccati}
For the linear-quadratic regulator problem
$\min\sum_{k=0}^{N-1}(x_k^\top Qx_k+u_k^\top R_wu_k)+x_N^\top
P_Nx_N$ s.t.\ $x_{k+1}=A_dx_k+B_du_k$, $Q\succeq0$, $R_w\succ0$, the
optimal cost-to-go is quadratic, $V_k(x)=x^\top P_kx$, with $P_k$
generated backward by the discrete Riccati recursion
\begin{equation}
P_k = Q + A_d^\top P_{k+1}A_d - A_d^\top P_{k+1}B_d(R_w+B_d^\top
P_{k+1}B_d)^{-1}B_d^\top P_{k+1}A_d, \qquad P_N \text{ given},
\label{eq:riccati-recursion}
\end{equation}
and the optimal control is the linear state feedback
$u_k^\star=-K_kx_k$, $K_k=(R_w+B_d^\top P_{k+1}B_d)^{-1}B_d^\top
P_{k+1}A_d$. If $(A_d,B_d)$ is stabilizable and $(A_d,Q^{1/2})$
detectable, $P_k\to P_\infty$ as the horizon recedes, where $P_\infty$
is the unique stabilizing solution of the discrete algebraic Riccati
equation (DARE), the fixed point of the Riccati recursion above.
\end{theorem}
\begin{proof}
By induction (Bellman's principle of optimality): assume $V_{k+1}(x)=x^\top
P_{k+1}x$. Then $V_k(x)=\min_u\left[x^\top Qx+u^\top R_wu+(A_dx+B_du)^\top
P_{k+1}(A_dx+B_du)\right]$ is an unconstrained convex quadratic in $u$
(since $R_w+B_d^\top P_{k+1}B_d\succ0$); setting its gradient to zero
gives $u^\star=-K_kx$ with $K_k$ as stated, and substituting back yields
$V_k(x)=x^\top P_kx$ with $P_k$ following the recursion above --
closing the induction. Convergence of the backward recursion to the
DARE's stabilizing fixed point under stabilizability/detectability is a
classical result; see Anderson \& Moore~\cite{andersonmoore}, Ch.~2, or
Bertsekas~\cite{bertsekasDP}, Vol.~I, Ch.~4.
\end{proof}

\begin{remark}
Theorem~\ref{thm:riccati} is restated here for self-containedness; it
matches the main course's own Theorem~15 (proved there by the identical
induction argument) and is the formal basis for
\S\ref{sec:mpc-theory}'s MPC-converges-to-LQR result.
\end{remark}

With this toolkit in hand, every later chapter's Theorem/Proof blocks
either specialize one of the results above or combine two or more of
them; cross-references back to this chapter are given at each use.

%=============================================================
\chapter{System Architecture and Data Schema}
\label{ch:architecture}
%=============================================================

\section{The schema layer}

Every other extension in this report ultimately reads or writes PSAT's
own \texttt{.con} device-data format. \pkg{schema/con\_\hspace{0pt}schema.json} was
built by reading PSAT's own \pkg{@XXclass/XXclass\hspace{0pt}.m} constructors
directly (never guessed), giving an authoritative column count
(\texttt{ncol}) for every device, and, for the highest-value devices this
project's labs use in detail, a traced column-by-column meaning obtained
by following the device's own \texttt{base.m}/\texttt{Fxcall.m} dynamics
code. Each entry carries an explicit confidence level:

\begin{itemize}
\item \textbf{verified} -- traced through source \emph{and}
  cross-checked against a known real value (e.g.\ the WSCC-9 case's
  \texttt{Syn.con} column 18 gives $H=23.64$\,s, the published textbook
  inertia constant). Covers \texttt{Bus}, \texttt{Line}, \texttt{Syn},
  \texttt{Exc}, \texttt{Tg} in full.
\item \textbf{data\_consistent} -- \texttt{ncol}/position confirmed from
  source; meaning inferred from position and plausible data, not traced
  through the differential-equation code. Covers \texttt{SW}, \texttt{PV},
  \texttt{PQ}, \texttt{Shunt}, \texttt{Fault}, \texttt{Breaker},
  \texttt{Supply}, \texttt{Demand}.
\item \textbf{unconfirmed} -- \texttt{ncol} authoritative; column meaning
  not resolved. Covers roughly twenty lower-priority devices (FACTS,
  HVDC, wind/DFIG, PSS, load-modelling variants).
\end{itemize}

\section{The Python I/O layer: \pkg{psat\_\hspace{0pt}io.py}}

\begin{figure}[H]
\centering
\includegraphics[width=0.52\textwidth]{Images/flowchart_psat_io.png}
\caption{Process flow of \pkg{psat\_io.py}.}
\end{figure}

\pkg{psat\_\hspace{0pt}io.py} reads a PSAT \texttt{.m} data file into named-column
\pkg{pandas} \texttt{DataFrame}s using the schema above, tested against
three real bundled files, including a re-verification of the WSCC-9
$H=23.64$\,s cross-check.

\section{The structured-results layer: \pkg{psat\_\hspace{0pt}results.m}}

\begin{figure}[H]
\centering
\includegraphics[width=0.46\textwidth]{Images/flowchart_psat_results.png}
\caption{Process flow of \pkg{psat\_results.m}.}
\end{figure}

\pkg{psat\_\hspace{0pt}results.m} collects PSAT's scattered post-run globals
(\texttt{Bus}, \texttt{Line}, \texttt{DAE}, \texttt{Syn}, \texttt{Varout},
\texttt{Settings}) into one struct. Two real bugs were found and fixed
while building it (full detail in Table~\ref{tab:bugs}): \texttt{DAE.y}
orders angles \emph{before} voltage magnitudes, the opposite of a first
draft's assumption, caught by checking that \texttt{DAE.y(10)} equalled
bus 1's known slack voltage of 1.04\,pu; and PSAT device objects are
old-style class instances, so \texttt{isfield(Syn,'con')} silently
returns \texttt{false} even when populated -- fixed with
\texttt{isempty(Syn.con)} instead.

\section{Parameter sweeps and physically-modified data files}

\pkg{psat\_\hspace{0pt}write\_\hspace{0pt}modified\_\hspace{0pt}datafile.m} writes a physically modified copy
of any device field (not only line status) to disk, because -- as
documented in Table~\ref{tab:bugs} -- editing PSAT's in-memory globals
and re-calling \texttt{runpsat('pf')} does \emph{not} take effect:
\texttt{runpsat} re-reads the original file from disk on every call
(\texttt{clpsat.readfile} defaults to 1). \pkg{psat\_\hspace{0pt}scenario\_\hspace{0pt}sweep.m}
generalizes this write/reload/collect pattern into a reusable sweep
driver, verified live on the WSCC-9 case: sweeping bus 6's real-power
demand from 0.5 to 2.1\,pu produces a clean, monotonically declining
voltage, $1.0183\to0.9838$\,pu.

\section{N-1 classification and secure ATC}

\begin{figure}[H]
\centering
\includegraphics[width=0.5\textwidth]{Images/flowchart_psat_n1_secure_atc.png}
\caption{Process flow of \pkg{psat\_n1\_classify.m} and \pkg{psat\_secure\_atc.m}.}
\end{figure}

\pkg{psat\_\hspace{0pt}n1\_\hspace{0pt}classify.m} runs N-1 screening and classifies each failure
as \emph{structural/islanding} (network disconnects, detected via
MATLAB's own \texttt{graph}/\texttt{conncomp}) versus
\emph{thermal/voltage} (network stays connected, a limit is exceeded).
Verified live on WSCC-9: exactly the three generator step-up transformer
lines (bus $1$--$4$, $2$--$7$, $3$--$9$) are flagged structural, the
remaining six are safe.

\pkg{psat\_\hspace{0pt}secure\_\hspace{0pt}atc.m} chains \texttt{cpfatc} across the full N-1
contingency list and returns the secure ATC, $\text{ATC}_{\text{secure}}
=\min(\text{ATC}_{\text{base}},\,\min_k \text{ATC}_k)$. Verified
live on \texttt{psat/tests/d\_003\_mdl.m}: base ATC = $989.974$, secure ATC
= $989.962$ (binding contingency: line 2). A separate, genuine
\emph{PSAT} bug (not this project's own code) was found and worked
around rather than silently patched: PSAT's internal N-1 sweep inside
\texttt{cpfatc} crashes on this MATLAB release once \texttt{Snapshot}
accumulates more than one element; \pkg{psat\_\hspace{0pt}secure\_\hspace{0pt}atc.m} exploits the
fact that \texttt{OPF.atc} is already set correctly before that internal
crash occurs.

\section{Linearization export for control design}

\begin{figure}[H]
\centering
\includegraphics[width=0.48\textwidth]{Images/flowchart_psat_export_ss.png}
\caption{Process flow of \pkg{psat\_export\_ss.m}.}
\end{figure}

\pkg{psat\_\hspace{0pt}export\_\hspace{0pt}ss.m} runs \texttt{pf}~$\to$~\texttt{sssa}~$\to$~
\texttt{fm\_abcd} and writes the resulting linearized $(A,B,C,D)$ to both
\texttt{.mat} and \texttt{.json}. Verified live on WSCC-9: $A$ is
$39\times39$, $B$ is $39\times6$ (3 AVR + 3 governor inputs), all
eigenvalues have non-positive real part. \pkg{psat\_\hspace{0pt}ss\_\hspace{0pt}loader.py} loads
the JSON export back into Python, completing the MATLAB$\to$Python
hand-off for \S\ref{ch:control}'s MPC/LQR work; a column-major-vs-row-major
flattening bug was found and fixed here (Table~\ref{tab:bugs}).

\section{Automated lab-report generation: \pkg{lab\_\hspace{0pt}report\_\hspace{0pt}generator.py}}

\begin{figure}[H]
\centering
\includegraphics[width=0.48\textwidth]{Images/flowchart_lab_report_generator.png}
\caption{Process flow of \pkg{lab\_report\_generator.py}.}
\end{figure}

\pkg{lab\_\hspace{0pt}report\_\hspace{0pt}generator.py} assembles a formatted, submittable
\texttt{.docx} lab report (title block, sections, tables, embedded
figures) from a run's own results, matching this course's own lab
``Report Requirements'' structure -- closing the gap left by PSAT's own
\texttt{pfrep}, which dumps plain text only. Verified two ways: every
generated document is read back with \pkg{python-docx} to confirm its
structure, and one test builds a real report embedding a real
\pkg{visualization.py} network plot, rendered all the way to PDF via
LibreOffice and visually inspected.

%=============================================================
\chapter{Deep-Integration Device Classes}
\label{ch:deviceclasses}
%=============================================================

Unlike every post-processing wrapper in Chapter~\ref{ch:architecture},
the two extensions in this chapter are genuine new PSAT \emph{device
types}, registered into PSAT's own class system so they participate
directly in \texttt{runpsat('pf')}/\texttt{('td')} exactly as
\texttt{Syn}/\texttt{Exc}/\texttt{Tg}/\texttt{PQ} do. Each required
additive edits (new lines only, nothing removed) to nine of PSAT's own
core dispatch files (\texttt{fm\_var.m}, \texttt{psat.m},
\texttt{fm\_call.m}, \texttt{fm\_dynidx.m}, \texttt{fm\_base.m},
\texttt{fm\_ncomp.m}, \texttt{fm\_inilf.m}, \texttt{fm\_restore.m},
\texttt{comp.ini}); backups of the pre-edit state of every file are kept
so either addition can be reverted independently.

\section{\pkg{@ESSclass}: a battery/storage device}

\begin{figure}[H]
\centering
\includegraphics[width=0.52\textwidth]{Images/flowchart_essclass.png}
\caption{Process flow of \pkg{@ESSclass}.}
\end{figure}

PSAT ships with no storage or battery device class of any kind
(confirmed by scanning every bundled \texttt{@XXclass} directory name and
\texttt{Contents.m} for ``storage'', ``battery'', ``ess''). \texttt{Ess}
closes this gap. \texttt{Ess.con} has 10 columns: \texttt{bus, Sn, Vn,
capacity[MWh], Pmax[pu], SoCmin, SoCmax, SoC0, Pref[pu], status}, with one
dynamic state per unit, the state of charge $\text{SoC}$ (as a fraction
of capacity):
\begin{equation}
\frac{d\,\text{SoC}}{dt} = -\frac{P_{\text{actual}}\cdot
S_{\text{mva}}}{\text{capacity}\cdot 3600}
\label{eq:ess}
\end{equation}
where $P_{\text{actual}}$ equals the commanded $P_{\text{ref}}$
(positive = discharging) unless the unit has saturated at
$\text{SoC}_{\min}$ or $\text{SoC}_{\max}$, in which case
$P_{\text{actual}}\to 0$ -- computed once in \texttt{ess\_effective\_
power.m} and shared by the network-injection equation
(\texttt{gcall.m}) and the state derivative (\texttt{fcall.m}) so the two
can never disagree about saturation. Known limitations, stated rather
than hidden: no reactive-power support ($Q=0$ always); $P_{\text{ref}}$
is a fixed constant, not a controllable time-varying dispatch signal
(a natural future extension); no internal resistance or thermal
derating.

\textbf{Verified live}, end-to-end against real PSAT 2.1.11 sessions,
with seven genuine bugs found by actually running it, including one
where the fix produced no errors and a clean exit but had, on direct
before/after numeric comparison, \emph{zero actual effect} -- caught only
by comparing numbers, not by trusting a clean run. Full defect log in
Table~\ref{tab:bugs}.

\section{\pkg{@GFMclass}: a grid-forming inverter (virtual synchronous machine)}

\begin{figure}[H]
\centering
\includegraphics[width=0.5\textwidth]{Images/flowchart_gfmclass.png}
\caption{Process flow of \pkg{@GFMclass}.}
\end{figure}

PSAT has no inverter-control device of any kind -- only the conventional
synchronous-machine model, \texttt{@SYclass}. \texttt{Gfm.con} has 9
columns: \texttt{bus, Sn, Vn, Hv[s], Dv[pu], Xv[pu], E[pu], Pref[pu],
status}, with two dynamic states per unit, virtual rotor angle
$\delta$ and virtual speed deviation $\omega$ (nominal 1\,pu), governed by
exactly the classical swing equation:
\begin{gather}
\frac{d\delta}{dt} = \text{rad}\cdot(\omega-1), \qquad \text{rad}=2\pi f_{\text{nominal}} \label{eq:gfm1}\\
\frac{d\omega}{dt} = \frac{P_{\text{ref}} - P_e - D_v(\omega-1)}{2H_v} \label{eq:gfm2}\\
P_e = \frac{EV}{X_v}\sin(\delta-\theta) \label{eq:gfm3}
\end{gather}
with $P_e$ (the electrical-power equation above) computed once in \texttt{gfm\_power.m} and
shared by \texttt{gcall}/\texttt{fcall}/\texttt{Gycall}/\texttt{Fxcall},
so all four always agree. The mismatch $P_{\text{ref}}-P_e$ excites the
swing dynamics -- the defining structural feature of a grid-\emph{forming}
(rather than grid-\emph{following}) inverter: its output is determined by
network conditions relative to its own internal state, not pinned to
$P_{\text{ref}}$ at every instant. This device also gives the main
course's Theorem 21 (``a virtual synchronous machine's closed-loop
frequency dynamics obey exactly the swing equation under substitution of
virtual for physical parameters'') a real PSAT implementation, not only a
standalone closed-form treatment.

\begin{theorem}[VSM small-signal dynamics and stability threshold, single-machine-infinite-bus case]
\label{thm:vsm}
Linearize the governing equations above about an equilibrium
$(\delta^\star,\omega=1)$ with the bus angle $\theta$ held fixed (the
single-machine-infinite-bus idealization), writing
$\Delta\delta=\delta-\delta^\star$, $\Delta\omega=\omega-1$. Then
\begin{equation}
\begin{bmatrix}\Delta\dot\delta\\\Delta\dot\omega\end{bmatrix}
= \underbrace{\begin{bmatrix}0 & \text{rad}\\ -\dfrac{K_v}{2H_v} & -\dfrac{D_v}{2H_v}\end{bmatrix}}_{A_{\text{vsm}}}
\begin{bmatrix}\Delta\delta\\\Delta\omega\end{bmatrix}, \qquad
K_v \;\triangleq\; \frac{EV}{X_v}\cos(\delta^\star-\theta),
\label{eq:vsm-linear}
\end{equation}
is \emph{algebraically identical} (same $2\times2$ block structure, same
roles for $H,D$, and the same synchronizing-coefficient definition
$K=\partial P_e/\partial\delta$ evaluated at $\delta^\star$) to the linearized
classical single-machine swing equation with physical inertia $H$ and
damping $D$ substituted by $H_v,D_v$. Consequently $A_{\text{vsm}}$'s
characteristic polynomial is
$2H_v\lambda^2+D_v\,\text{rad}^{-1}\lambda\cdot\text{rad}+K_v\,\text{rad}=0$,
i.e.\ (absorbing $\text{rad}$ into a time rescaling, the standard
per-unit-radian swing-equation convention)
$2H_v\lambda^2+D_v\lambda+K_v=0$, with roots
\begin{equation}
\lambda_{1,2} = \frac{-D_v\pm\sqrt{D_v^2-8H_vK_v}}{4H_v}.
\label{eq:vsm-eigs}
\end{equation}
The equilibrium is locally asymptotically stable if and only if
$D_v>0$ and $K_v>0$ (Routh--Hurwitz necessary and sufficient conditions
for a second-order real polynomial); since $K_v$ changes sign exactly
when $\delta^\star-\theta$ crosses $90^\circ$, stability is lost
\emph{discontinuously} at that angle, not gradually.
\end{theorem}
\begin{proof}
Differentiate the governing swing equations above and evaluate the
Jacobian at $(\delta^\star,1)$: $\partial(\Delta\dot\delta)/
\partial\Delta\delta=0$, $\partial(\Delta\dot\delta)/\partial\Delta
\omega=\text{rad}$ directly from the angle equation; from the speed
and power equations, $\partial(\Delta\dot\omega)/
\partial\Delta\delta=-\tfrac1{2H_v}\partial P_e/\partial\delta=
-K_v/(2H_v)$ and $\partial(\Delta\dot\omega)/\partial\Delta\omega=
-D_v/(2H_v)$, giving the linearized state matrix above. The characteristic
polynomial of a companion-form $2\times2$ matrix
$\begin{bmatrix}0&a\\b&c\end{bmatrix}$ is $\lambda^2-c\lambda-ab=0$;
substituting $a=\text{rad}$, $b=-K_v/(2H_v)$, $c=-D_v/(2H_v)$ and
rescaling time by $\text{rad}$ (equivalent to measuring $\lambda$ in
electrical radians/s instead of per-unit speed units, the standard swing-
equation normalization -- see Kundur~\cite{kundur}, Ch.~11, or Machowski,
Bialek \& Bumby~\cite{machowski}, Ch.~5) yields the eigenvalue formula
above. Both roots have negative real part iff the
polynomial's two coefficients $D_v,K_v$ (with leading coefficient
$2H_v>0$) are both positive -- the elementary 2-term Routh array test
for a quadratic, equivalent to Vieta's formulas giving
$\lambda_1\lambda_2=K_v/(2H_v)>0$ and $\lambda_1+\lambda_2=-D_v/(2H_v)<0$
simultaneously forcing both roots into the open left half-plane.
\end{proof}

\begin{corollary}[Pinned decay rate under uniform parameterization]
\label{cor:pinned-decay}
If $D_v^2<8H_vK_v$ (the underdamped case), the eigenvalue formula above gives a
complex-conjugate pair with real part exactly $-D_v/(4H_v)$,
\emph{independent of $K_v$'s magnitude} -- only the oscillation
frequency (the imaginary part) depends on $K_v$. This is the closed-form
origin of this report's \S\ref{sec:stability-opf-theory} empirical
finding that a uniformly $H$/$D$-parameterized multi-machine system's
dominant eigenvalue real part stays pinned at $-D/(4H)$ for every stable
configuration.
\end{corollary}

\textbf{Verified live}; building it caused a real small-signal
instability to be discovered (eigenvalue check
$\text{eig}(\text{DAE.Fx}-\text{DAE.Fy}(\text{DAE.Gy}\backslash
\text{DAE.Gx}))$ showing a genuinely unstable mode at $+25.9$), traced
to a sign-convention error in how generation is injected into PSAT's own
power-balance equation -- and the same, previously undetected bug was
then found and fixed in the already-shipped \texttt{@ESSclass} too. Full
detail in Table~\ref{tab:bugs}.

%=============================================================
\chapter{Distribution, Fault, and Multi-Energy Network Analysis}
\label{ch:distribution}
%=============================================================

\section{Three-phase unbalanced distribution power flow}

\begin{figure}[H]
\centering
\includegraphics[width=0.5\textwidth]{Images/flowchart_three_phase_distribution.png}
\caption{Process flow of \pkg{three\_phase\_distribution.py}.}
\end{figure}

PSAT is a positive-sequence-only transmission tool;
\pkg{three\_\hspace{0pt}phase\_\hspace{0pt}distribution\hspace{0pt}.py}
fills the unbalanced-distribution gap via the standard
forward--backward sweep (FBS) method on a radial feeder, each line
segment described by a $3\times3$ phase impedance matrix $Z_{\text{line}}$
(allowing mutual phase coupling, which a positive-sequence model cannot
represent at all):
\begin{gather}
\text{backward:} \quad  I_{\text{line}} = \textstyle\sum_{\downarrow\text{loads}} I_{\text{load}}(V) + \sum_{\downarrow\text{lines}} I_{\text{line}} \\
\text{forward:} \quad  V_{\text{to}} = V_{\text{from}} - Z_{\text{line}}\,I_{\text{line}}
\end{gather}
(sums taken over the loads and line segments electrically downstream of
the current bus)
iterated to convergence. Verified against the published IEEE 4-bus test
feeder (balanced-load variant, giving an independent cross-check against
PYPOWER on the same topology) and a hand-computable 2-bus unbalanced
case. A real, consequential bug (Table~\ref{tab:bugs}) was found in the
MATLAB port: \texttt{'} is MATLAB's \emph{conjugate} transpose for a
complex array, silently producing a fully converged but wrong
phase-B/phase-C-swapped answer.

\begin{proposition}[FBS as a fixed-point iteration]
\label{prop:fbs}
Write one backward--forward sweep as $V^{(k+1)}=T(V^{(k)})$, where $T$
composes the backward current calculation $I_{\text{line}}(V)$
(constant-power loads give $I_{\text{load},\phi}=
(S_\phi/V_\phi)^{\star}$, a function of $V$) with the forward voltage
update $V_{\text{to}}=V_{\text{from}}-Z_{\text{line}}I_{\text{line}}$. If, on a neighborhood $\mathcal{V}$ of the true solution
$V^\star$, $T$ is Lipschitz with constant $\gamma<1$ in $V$ (a
sufficient condition is that the per-unit voltage drop across every
line segment, $\|Z_{\text{line}}I_{\text{line}}\|/\|V\|$, stays
small enough that the constant-power load's current-vs-voltage
sensitivity does not dominate the line-drop sensitivity --
quantitatively, $\gamma\approx\max_\phi\left|Z_{\text{line}}S_\phi/
V_\phi^2\right|$, the standard FBS convergence-radius estimate; see Kersting,
\emph{Distribution System Modeling and Analysis}~\cite{kersting},
Ch.~10, and Cheng \& Shirmohammadi~\cite{chengshirmohammadi}), then by
Theorem~\ref{thm:banach} the FBS iteration converges to the unique fixed
point $V^\star\in\mathcal V$ geometrically. Lightly loaded radial
feeders (the regime every IEEE test feeder and this module's own test
cases operate in) satisfy this condition comfortably; heavily loaded or
very long feeders can violate it, the standard practical failure mode of
plain FBS reported throughout the distribution-analysis literature.
\end{proposition}

\section{Short-circuit fault analysis}

\begin{figure}[H]
\centering
\includegraphics[width=0.5\textwidth]{Images/flowchart_short_circuit.png}
\caption{Process flow of \pkg{short\_circuit.py}.}
\end{figure}

\pkg{short\_\hspace{0pt}circuit.py} computes symmetrical (three-phase) and
unsymmetrical (SLG, LL, DLG) fault currents via the $Z_{\text{bus}}$
method and symmetrical components, the standard protection-and-
coordination calculation PSAT's own \texttt{Fault.con} (a pure
time-domain disturbance device) cannot perform:
\begin{gather}
I_f^{3\phi} = \frac{E_a}{Z_{1,kk}+Z_f} \\
I_f^{\text{SLG}} = \frac{3E_a}{Z_{0,kk}+Z_{1,kk}+Z_{2,kk}+3Z_f} \\
I_f^{\text{LL}} = \frac{E_a}{Z_{1,kk}+Z_{2,kk}+Z_f}
\end{gather}
with the DLG case solved via the standard parallel $Z_1/Z_2/Z_0$
combination. Verified against PYPOWER's \texttt{case9} positive-sequence
$Y_{\text{bus}}$ and a closed-form single-generator infinite-bus
calculation. The MATLAB port (\pkg{psat\_\hspace{0pt}short\_\hspace{0pt}circuit.m}) was
cross-validated against an independent Python run reading the identical
PSAT file, matching $Z_{kk}=0.0007+j0.0514$, $\left|I_f\right|=19.4644$\,pu in both
languages to the digit.

\begin{theorem}[Fault current from the Thevenin/superposition argument]
\label{thm:zbus-fault}
For a bolted three-phase fault at bus $k$ of a linear network with
pre-fault internal EMF $E_a$ behind positive-sequence bus impedance
matrix $Z_1=Y_{\text{bus},1}^{-1}$, the fault current is
$I_f^{3\phi}=E_a/(Z_{1,kk}+Z_f)$, and the fault-induced voltage change
at any bus $m$ is $\Delta V_m=-Z_{1,mk}I_f^{3\phi}$.
\end{theorem}
\begin{proof}
By superposition, the faulted network's response equals the pre-fault
response plus the response to a single current source of value
$-I_f^{3\phi}$ injected at bus $k$ alone (Thevenin's theorem applied at
bus $k$: the fault is modeled as closing a switch to a fault impedance
$Z_f$, equivalent to injecting the current that switch carries while
every other source is deactivated for the purposes of computing the
\emph{change} in network state). The bus-impedance matrix $Z_1=
Y_{\text{bus},1}^{-1}$ exists because $Y_{\text{bus},1}$ is the
network's own admittance Laplacian plus shunt generator admittances --
nonsingular whenever at least one generator is connected, by the same
argument as Theorem~\ref{thm:laplacian} applied to the augmented
(grounded) network. The change in bus $m$'s voltage due to a unit
current injected at bus $k$ alone is, by definition of the inverse,
exactly column $k$ of $Z_1$, i.e.\ $\Delta V_m=-Z_{1,mk}I_f^{3\phi}$ for
an injected current of $-I_f^{3\phi}$. At the fault bus itself, the
boundary condition $V_k=Z_fI_f^{3\phi}$ (Ohm's law through the fault
impedance) combined with $V_k=E_a-Z_{1,kk}I_f^{3\phi}$ (pre-fault voltage
minus the self-impedance drop) gives $E_a-Z_{1,kk}I_f^{3\phi}=
Z_fI_f^{3\phi}$, solved directly for $I_f^{3\phi}=E_a/(Z_{1,kk}+Z_f)$.
The unsymmetrical-fault formulas follow the identical argument applied
in the decoupled positive/negative/zero-sequence networks, connected at
the fault point according to the fault type's own boundary conditions
(Fortescue's symmetrical-component transformation block-diagonalizes a
balanced three-phase network into three independent single-phase
sequence networks); see Glover, Overbye \&
Sarma~\cite{glover}, Ch.~9--10, or Grainger \&
Stevenson~\cite{graingerstevenson}, Ch.~11--12, for the full derivation
of each fault type's sequence-network interconnection.
\end{proof}

\section{Gas-pipeline networks: the Weymouth equation}

\begin{figure}[H]
\centering
\includegraphics[width=0.5\textwidth]{Images/flowchart_gas_network.png}
\caption{Process flow of \pkg{gas\_network.py}.}
\end{figure}

\pkg{gas\_\hspace{0pt}network.py} deliberately mirrors the structure of this course's
own Newton--Raphson power-flow treatment, substituting pressure for
voltage. For a pipe $(i,j)$ with Weymouth constant $K_{ij}>0$:
\begin{equation}
q_{ij} = \text{sign}(p_i^2-p_j^2)\sqrt{K_{ij}\left|p_i^2-p_j^2\right|}
\label{eq:weymouth}
\end{equation}
exactly analogous to DC power flow's $P_{ij}=B_{ij}(\theta_i-\theta_j)$,
except nonlinear in pressure (gas is compressible) rather than linear in
angle. Nodal balance at every non-reference node,
$\text{supply}_i-\text{demand}_i-\sum_j q_{ij}=0$, is solved by
Newton--Raphson on the unknown pressures directly (not pressure-squared,
to keep positivity visible during iteration) with a finite-difference
Jacobian, cross-validated against \texttt{scipy.optimize.fsolve} on every
test case. The MATLAB port reproduces the Python/\texttt{fsolve}-
cross-validated 3-node result ($C=60.18$, $D=50.29$) to 2 decimal places,
after a real bug (nested MATLAB functions closing over outer data
silently returning \texttt{NaN}) was found and fixed by switching to
ordinary local functions with explicit arguments.

\begin{theorem}[Local quadratic convergence of Newton--Raphson]
\label{thm:newton}
Let $F:\mathbb{R}^n\to\mathbb{R}^n$ be continuously differentiable with
Jacobian $J(x)$ Lipschitz continuous with constant $L$ near a solution
$x^\star$ ($F(x^\star)=0$), and suppose $J(x^\star)$ is nonsingular. The
Newton iteration $x_{k+1}=x_k-J(x_k)^{-1}F(x_k)$ converges
\emph{quadratically} for any $x_0$ sufficiently close to $x^\star$:
$\|x_{k+1}-x^\star\|\le c\|x_k-x^\star\|^2$ for a constant $c$ depending
on $L$ and $\|J(x^\star)^{-1}\|$.
\end{theorem}
\begin{proof}
By Taylor's theorem with integral remainder,
$0=F(x^\star)=F(x_k)+J(x_k)(x^\star-x_k)+r_k$ with
$\|r_k\|\le\tfrac{L}{2}\|x^\star-x_k\|^2$ (the standard Lipschitz-
Jacobian remainder bound). Rearranging,
$x^\star-x_k+J(x_k)^{-1}F(x_k)=-J(x_k)^{-1}r_k$, i.e.\
$x_{k+1}-x^\star=J(x_k)^{-1}r_k$, so $\|x_{k+1}-x^\star\|\le
\|J(x_k)^{-1}\|\cdot\tfrac{L}{2}\|x_k-x^\star\|^2$. Continuity of
$J(\cdot)^{-1}$ near $x^\star$ (from $J(x^\star)$ nonsingular) bounds
$\|J(x_k)^{-1}\|$ by a constant on a small enough neighborhood, giving
the stated quadratic bound; see Nocedal \&
Wright~\cite{nocedalwright}, Thm.~11.2, or Ortega \& Rheinboldt. This is
restated here from the main course's own identical Theorem~2 (there
proved for the AC power-flow mismatch function) because
\pkg{gas\_\hspace{0pt}network.py} solves the structurally analogous nodal-balance
root-finding problem $F(p)=\text{supply}-\text{demand}-\sum q(p)=0$ on
the unknown nodal pressures, with the finite-difference Jacobian
approximating $J(p)$ to the accuracy needed for this theorem's
hypotheses to hold in the discretized implementation.
\end{proof}

\begin{remark}
Unlike the DC power-flow map (linear, Theorem~\ref{thm:laplacian}
guarantees a unique solution unconditionally once the reference node is
fixed), the Weymouth nodal-balance system is
genuinely nonlinear in pressure, so this report does not claim a
global existence/uniqueness theorem for it -- only the local
convergence guarantee of Theorem~\ref{thm:newton} once Newton's method
is already close to a solution, consistent with every test case in this
module being cross-validated against \texttt{scipy.optimize.fsolve}
from a physically reasonable flat start (uniform nominal pressure)
rather than an arbitrary one.
\end{remark}

\section{Network visualization}

\begin{figure}[H]
\centering
\includegraphics[width=0.46\textwidth]{Images/flowchart_visualization.png}
\caption{Process flow of \pkg{visualization.py}.}
\end{figure}

\pkg{visualization\hspace{0pt}.py} draws schematic network diagrams (via
\texttt{networkx} layout, since PSAT data files carry no bus
coordinates) with optional voltage-magnitude node coloring and
line-loading edge coloring, verified against the WSCC-9 structural-failure
lines Chapter~\ref{ch:architecture}'s N-1 classifier already confirmed
live. The MATLAB port, \pkg{psat\_\hspace{0pt}visualize.m}, plots PSAT's own
currently-loaded network directly (never a separately re-typed topology).

%=============================================================
\chapter{Optimization and Economic Dispatch Extensions}
\label{ch:optimization}
%=============================================================

PSAT's \texttt{opf} is a single, centralized, deterministic solve. This
chapter's extensions add distributed optimization, uncertainty (both
stochastic and robust), storage dispatch, market economics, multi-period
scheduling, and dynamic-stability-constrained dispatch -- none of which
PSAT has any native capability for.

\section{Distributed optimization: dual decomposition, ADMM, and consensus}
\label{sec:decomp-theory}

\begin{figure}[H]
\centering
\includegraphics[width=0.5\textwidth]{Images/flowchart_decentralized_opt.png}
\caption{Process flow of \pkg{decentralized\_opt.py}.}
\end{figure}

\pkg{decentralized\_\hspace{0pt}opt.py} implements three structurally distinct
distributed-dispatch algorithms, each benchmarked against the closed-form
centralized equal-incremental-cost optimum,
\begin{equation}
\lambda^\star = \frac{d+\sum_i b_i/(2a_i)}{\sum_i 1/(2a_i)}, \qquad
P_i^\star = \frac{\lambda^\star-b_i}{2a_i}.
\label{eq:centralized-ed}
\end{equation}

\textbf{Dual decomposition} (separable sum-coupling constraint
$\sum_i P_i=d$): each generator independently minimizes its own cost
minus $\lambda P_i$; the scalar price updates by subgradient ascent,
\begin{equation}
\lambda^{(k+1)} = \lambda^{(k)} + \alpha\left(d-\textstyle\sum_i P_i^{(k)}\right),
\qquad \alpha = \left(\textstyle\sum_i \tfrac{1}{2a_i}\right)^{-1}.
\label{eq:dual-decomp}
\end{equation}
Verified on a 2-generator test system ($a=[0.04,0.05]$, $b=[9.6,11.2]$,
$d=400$): both the closed-form and the iterative solver converge to
$P=[231.11,168.89]$, $\lambda=28.0889$, in 3 iterations.

\textbf{ADMM consensus on a shared variable} (the tie-line flow
$P_{\text{tie}}$, which dual decomposition cannot handle because there
is no separable sum-coupling constraint -- both areas must agree on the
\emph{same} value): each area solves a locally penalized sub-problem,
\begin{equation}
t_a \leftarrow \arg\min_{t\in[-\bar t,\bar t]}\; \text{cost}_a\left(d_a+t\right) + \tfrac{\rho}{2}(t-(z-u_a))^2,
\label{eq:admm}
\end{equation}
followed by the standard two-block ADMM update
$z\leftarrow\tfrac12(t_a+t_b)$, $u_a\leftarrow u_a+(t_a-z)$,
$u_b\leftarrow u_b+(t_b-z)$. Verified on a 2-area test system: converges
to $P_{\text{tie}}=61.1111$, $P_a=211.1111$, $P_b=188.8889$.

\textbf{Metropolis--Hastings consensus dispatch} (fully decentralized,
no coordinator, neighbor-to-neighbor communication over an adjacency
graph), using symmetric doubly-stochastic weights
$W_{ij}=1/(1+\max(\deg_i,\deg_j))$ for $i\ne j$, $W_{ii}=1-\sum_{j\ne i}W_{ij}$
-- required to be symmetric because a row-stochastic-only matrix was
found, in this project's own development, to converge to the wrong,
degree-weighted fixed point rather than the true consensus average.

\begin{theorem}[Strong duality and subgradient convergence for separable economic dispatch]
\label{thm:dual-decomp}
The quadratic economic-dispatch problem $\min_P\sum_ia_iP_i^2+b_iP_i$
s.t.\ $\sum_iP_i=d$, $P_i\ge0$ is convex with an affine equality
constraint, so Slater's condition (Assumption~\ref{as:slater}) holds
trivially (any feasible point is already ``strict'' for a problem with
no inequality constraints beyond box bounds met with slack) and
Theorem~\ref{thm:strong-duality} applies: strong duality holds and the
KKT stationarity condition $2a_iP_i+b_i-\lambda=0$ (for $P_i$ in the
interior of its bound) gives exactly the closed form stated above
for the centralized optimum, with $\lambda^\star$ the shared dual price.
Moreover, the subgradient iteration given above converges to
$\lambda^\star$ at the sublinear rate $O(1/\sqrt{k})$ for a general
convex dual function and \emph{geometrically} here, because the dual
function $d(\lambda)=\sum_i\min_{P_i\ge0}(a_iP_i^2+b_iP_i-\lambda
P_i)+\lambda d$ is itself twice differentiable and strongly
concave in $\lambda$ with curvature $-\sum_i\tfrac1{2a_i}$ on the
unconstrained-$P_i$ region, making that same subgradient iteration exactly
Newton's method on $d(\lambda)$ with the problem's own curvature as the
step size -- the formal reason the module converges to machine precision
in exactly 3 iterations on its 2-generator reference case rather than the
many iterations a generic (curvature-oblivious) subgradient step would
need. See Bertsekas~\cite{bertsekasnlp}, \S6.1, or Boyd \&
Vandenberghe~\cite{boydconvex}, \S5.5.3, for the general dual-
decomposition convergence theory.
\end{theorem}

\begin{theorem}[ADMM convergence]
\label{thm:admm}
For $\min_{x,z}f(x)+g(z)$ s.t.\ $Ax+Bz=c$ with $f,g$ closed, proper,
convex, the ADMM iterates
$x^{k+1}=\arg\min_xf(x)+\tfrac\rho2\|Ax+Bz^k-c+u^k\|_2^2$,
$z^{k+1}=\arg\min_zg(z)+\tfrac\rho2\|Ax^{k+1}+Bz-c+u^k\|_2^2$,
$u^{k+1}=u^k+Ax^{k+1}+Bz^{k+1}-c$ satisfy $f(x^k)+g(z^k)\to p^\star$
(the optimal value) and the primal residual $Ax^k+Bz^k-c\to0$, for any
$\rho>0$, provided a saddle point of the unaugmented Lagrangian exists.
\end{theorem}
\begin{proof}[Proof sketch]
Boyd, Parikh, Chu, Peleato \& Eckstein~\cite{boyd}, \S3.2, prove this
via a Lyapunov-function argument: $V^k=\tfrac1\rho\|u^k-u^\star\|_2^2+
\rho\|B(z^k-z^\star)\|_2^2$ is shown to be non-increasing along the
ADMM iterates (using convexity of $f,g$ and the optimality conditions at
a saddle point $(x^\star,z^\star,u^\star)$), and the decrease
$V^k-V^{k+1}$ dominates the sum of squared primal and dual residual
norms, which therefore must vanish; convergence of the objective then
follows from convexity and the residuals vanishing. Applied to
the ADMM local update given above with $x=(t_a,P_a)$, $z=(t_b,P_b)$, $A$/$B$ encoding the
consensus constraint $t_a=t_b$: this is exactly the two-block ADMM
pattern, so Theorem~\ref{thm:admm} guarantees the module's observed
convergence to the shared tie flow $P_{\text{tie}}=61.1111$ is not
coincidental -- it is the generic behavior ADMM theory predicts for
\emph{any} convex two-area split of this problem, for any $\rho>0$
(though the convergence \emph{rate} does depend on $\rho$, which is why
this project's own earlier work needed to move from a coarse grid search
to a proper bounded 1-D solver for the ADMM inner problem --
Table~\ref{tab:bugs} -- to realize the rate the theorem's hypotheses
actually promise).
\end{proof}

\begin{corollary}[Metropolis--Hastings consensus dispatch converges to the true average]
\label{cor:mh-dispatch}
The weight matrix $W$ defined above is symmetric by construction
($W_{ij}=W_{ji}$) and doubly stochastic (row sums
$\sum_{j\ne i}W_{ij}+W_{ii}=1$ by construction, and column sums equal row
sums by symmetry), and irreducible whenever the communication graph is
connected. Theorem~\ref{thm:doubly-stochastic} therefore guarantees each
agent's local price estimate $\lambda_i^{(k)}$ converges geometrically to
the true network-wide average price, \emph{not} a degree-weighted
average -- the exact property a naive row-stochastic-only $W$ lacks, and
the root cause of the bug this project's own Week~12 development and
this module's own first draft both independently hit (Table~\ref{tab:bugs}).
\end{corollary}

\section{Stochastic and robust dispatch under uncertainty}
\label{sec:chance-theory}

\begin{figure}[H]
\centering
\includegraphics[width=0.52\textwidth]{Images/flowchart_stochastic_opf.png}
\caption{Process flow of \pkg{stochastic\_opf.py}.}
\end{figure}

\textbf{Two-stage stochastic recourse} (\pkg{stochastic\_\hspace{0pt}opf.py}):
first-stage (here-and-now) commitment, second-stage (wait-and-see)
recourse per scenario $s$ with probability $\pi_s$:
\begin{equation}
\min_{x}\; c^\top x + \sum_s \pi_s\, Q_s(x), \qquad
Q_s(x)=\min_{y_s\ge0}\; q^\top y_s \;\; \text{s.t. recourse feasibility}.
\end{equation}
Verified to reproduce Value of the Stochastic Solution
$\text{VSS}=\$138.0$ on this project's reference scenario, and
cross-checked against brute-force enumeration over a discretized
first-stage grid.

\begin{theorem}[Ordering of the stochastic-programming bounds, and $\text{VSS}\ge0$]
\label{thm:vss}
Define the recourse problem value
$\text{RP}=\min_x c^\top x+\mathbb E_\xi[Q(x,\xi)]$, the wait-and-see
value $\text{WS}=\mathbb E_\xi[\min_x c^\top x+Q(x,\xi)]$, the
expected-value solution $\bar x=\arg\min_x c^\top x+Q(x,\mathbb
E[\xi])$, and $\text{EEV}=c^\top\bar x+\mathbb E_\xi[Q(\bar x,\xi)]$
(the true expected cost of deploying the EV-problem's solution under
the real uncertainty). If $Q(x,\cdot)$ is convex in $\xi$ for each fixed
$x$ (true whenever $Q$ is itself the value function of a convex
recourse LP/QP, as it is here), then
\begin{equation}
\text{WS}\;\le\;\text{RP}\;\le\;\text{EEV}, \qquad
\text{VSS}\triangleq\text{EEV}-\text{RP}\;\ge\;0.
\label{eq:vss-chain}
\end{equation}
\end{theorem}
\begin{proof}
$\text{RP}\le\text{EEV}$ immediately: $\text{RP}$ minimizes over
\emph{all} $x$ while $\text{EEV}$ evaluates the (generally suboptimal)
fixed choice $\bar x$, so $\text{RP}=\min_xc^\top x+\mathbb
E_\xi[Q(x,\xi)]\le c^\top\bar x+\mathbb E_\xi[Q(\bar x,\xi)]=
\text{EEV}$. For $\text{WS}\le\text{RP}$: for every fixed $\xi$,
$\min_xc^\top x+Q(x,\xi)\le c^\top x+Q(x,\xi)$ for the particular $x$
that minimizes the expectation, so taking expectations of the
left-hand side's pointwise minimum can only decrease the value relative
to fixing one $x$ across all $\xi$ -- formally, $\mathbb
E_\xi\left[\min_xc^\top x+Q(x,\xi)\right]\le\min_x\mathbb E_\xi\left[c^\top
x+Q(x,\xi)\right]$ because the expectation of a pointwise minimum never
exceeds the minimum of the expectation (interchanging $\min$ and
$\mathbb E$ can only help the wait-and-see planner, who chooses $x$
\emph{after} seeing $\xi$). See Birge \&
Louveaux~\cite{birgelouveaux}, \S1.2--1.3, for the complete
derivation and the role of convexity. The module's own verified
$\text{VSS}=\$138.0>0$ on its reference scenario is therefore not a
numerical coincidence but a guaranteed structural property of any
convex two-stage recourse problem with genuinely scenario-dependent
recourse costs.
\end{proof}

\textbf{Gaussian chance-constrained reserve sizing}: the smallest reserve
$R^\star$ such that $\Pr(\text{shortfall}>R)\le\epsilon$ under a Gaussian
forecast-error model is
\begin{equation}
R^\star = \sigma\,\Phi^{-1}(1-\epsilon),
\label{eq:chance}
\end{equation}
reproducing $R_{90}=2.957$, $R_{95}=3.795$, $R_{99}=5.368$\,MW at
$\sigma=2.3074$\,MW.

\begin{theorem}[The Gaussian quantile reserve is the tightest valid reserve]
\label{thm:chance-tight}
Let shortfall $\sim\mathcal N(0,\sigma^2)$. Among all $R\ge0$ satisfying
$\Pr(\text{shortfall}>R)\le\epsilon$, $R^\star=\sigma\Phi^{-1}(1-\epsilon)$
is the smallest.
\end{theorem}
\begin{proof}
$\Pr(\text{shortfall}>R)=1-\Phi(R/\sigma)$ is a strictly decreasing,
continuous function of $R$ (since $\Phi$ is a strictly increasing CDF
with a everywhere-positive Gaussian density). The constraint
$1-\Phi(R/\sigma)\le\epsilon$ is therefore equivalent to
$R\ge\sigma\Phi^{-1}(1-\epsilon)$ (applying the strictly increasing
function $\Phi^{-1}$ to both sides of $\Phi(R/\sigma)\ge1-\epsilon$),
and the smallest $R$ satisfying a $\ge$ constraint is attained with
equality, $R^\star=\sigma\Phi^{-1}(1-\epsilon)$. Strict monotonicity
(not just monotonicity) is what rules out a smaller $R$ also
satisfying the constraint, and is exactly the property a non-Gaussian,
possibly non-strictly-monotonic error CDF can lack, which is why the
main course's own Theorem~13 (and this report's
\S\ref{ch:limitations}) flags the Gaussian assumption explicitly as a
modeling choice, not a universal law.
\end{proof}

\begin{figure}[H]
\centering
\includegraphics[width=0.5\textwidth]{Images/flowchart_robust_opf.png}
\caption{Process flow of \pkg{robust\_opf.py}.}
\end{figure}

\begin{theorem}[Bertsimas--Sim robust counterpart]
\label{thm:bertsimas-sim}
Consider an uncertain linear constraint $\sum_ia_ix_i\le b$ where each
coefficient $a_i=\bar a_i+\xi_i\hat a_i$ may deviate from its nominal
value $\bar a_i$ by up to $\hat a_i\ge0$, with $\xi_i\in[-1,1]$, and at
most $\Gamma\in[0,n]$ of the $n$ uncertain coefficients are allowed to
simultaneously reach their bound (the ``budget of uncertainty''). The
robust requirement ``the constraint holds for every realization
consistent with the budget'' is equivalent to the single deterministic
constraint
\begin{equation}
\sum_i\bar a_ix_i \;+\; \max_{\{S\cup\{t\}:\,S\subseteq\{1,\ldots,n\},\left|S\right|\le\lfloor\Gamma\rfloor,\,t\notin S\}}
\left[\sum_{i\in S}\hat a_i\left|x_i\right| + (\Gamma-\lfloor\Gamma\rfloor)\hat a_t\left|x_t\right|\right] \;\le\; b,
\label{eq:bs-robust-raw}
\end{equation}
and this inner maximization itself admits the exact linear-programming
dual reformulation
\begin{equation}
\sum_i\bar a_ix_i + \Gamma z + \sum_ip_i \le b, \quad
z+p_i\ge\hat a_i\left|x_i\right|\;\;\forall i, \quad z,p_i\ge0,
\label{eq:bs-robust-lp}
\end{equation}
with $(z,p_i)$ new decision variables -- a tractable LP/QP of the same
order as the nominal problem, with no explicit enumeration of worst-case
scenarios required.
\end{theorem}
\begin{proof}[Proof sketch]
The raw combinatorial inner maximization above is a linear program
in the deviation pattern (which coefficients deviate, and whether a
single fractional-budget coefficient partially deviates): choosing the
$\Gamma$ largest-impact terms $\hat a_i\left|x_i\right|$ to set at their bound is
optimal by a standard exchange argument (swapping any smaller-impact
chosen deviation for a larger-impact unchosen one cannot decrease the
sum), which is exactly what the combinatorial maximization expresses.
Writing that LP's own dual (a textbook LP-duality application,
Theorem~\ref{thm:strong-duality} specialized to this particular
bounded-knapsack-like LP) replaces the combinatorial $\max$ over
subsets with the continuous variables $(z,p_i)$ in the LP reformulation
above, strong duality guaranteeing the two are
equal; see Bertsimas \& Sim~\cite{bertsimas}, \S2--3, for the full
derivation, including the case of simultaneously uncertain
objective coefficients (used identically here for uncertain demand
entering the balance constraint's right-hand side, a special case with
$x_i\equiv1$ coefficients on the uncertain terms). This is the formal
reason \pkg{robust\_\hspace{0pt}opf.py} never needs to solve an explicit inner
maximization or enumerate worst cases: the LP reformulation above is
solved directly as part of the outer dispatch QP.
\end{proof}

\begin{corollary}[Monotonicity of the price of robustness]
\label{cor:price-robustness}
The optimal value of the LP reformulation above is non-decreasing in
$\Gamma$ for $\Gamma\in[0,n]$, since enlarging $\Gamma$ only enlarges the
feasible region of the raw combinatorial inner maximization
(more simultaneous worst-case deviations are permitted), which can only
tighten -- never loosen -- the resulting deterministic constraint
on the outer minimization. This is the formal
reason this module's own ``price of robustness'' sweep is monotonic by
construction, not merely an empirical regularity: $\Gamma=0$ recovers
the nominal (deterministic) problem exactly, and $\Gamma=n$ recovers the
fully conservative box-uncertainty worst case.
\end{corollary}

\section{Battery/storage arbitrage dispatch}

\begin{figure}[H]
\centering
\includegraphics[width=0.46\textwidth]{Images/flowchart_battery_storage.png}
\caption{Process flow of \pkg{battery\_storage.py}.}
\end{figure}

\pkg{battery\_\hspace{0pt}storage.py} solves LP-optimal price-arbitrage dispatch,
\begin{equation}
\min_{P_t}\; \sum_t \pi_t P_t \quad \text{s.t.}\quad
\text{SoC}_{t+1}=\text{SoC}_t+P_t\,\Delta t,\;\;
P_t\in[-\bar P,\bar P],\;\; \text{SoC}_t\in[\text{SoC}_{\min},\text{SoC}_{\max}],
\label{eq:battery-lp}
\end{equation}
via \pkg{cvxpy} (Python) / \texttt{linprog} (MATLAB). Verified on a
10\,kWh/3\,kW battery against a 12-period price series: both languages
reproduce the identical bang-bang schedule
$P=[2,\allowbreak-3,\allowbreak3,\allowbreak3,\allowbreak-3,\allowbreak-3,\allowbreak-3,\allowbreak3,\allowbreak3,\allowbreak-1,\allowbreak-3,\allowbreak-3]$\,kW, cost = $-\$237.0$, to the exact
value -- a full, no-discrepancy match, with no bugs found in either
implementation on the first live run.

\begin{proposition}[Bang-bang structure of the arbitrage optimum]
\label{prop:bangbang}
The battery arbitrage dispatch above is a linear program (linear objective, linear
equality/inequality constraints), so by the fundamental theorem of
linear programming an optimal solution is attained at a vertex (basic
feasible solution) of the feasible polytope -- see Bertsimas \&
Tsitsiklis~\cite{bertsimastsitsiklis}, Thm.~2.7 (restated as the main
course's own Theorem~12). A vertex of this particular polytope has every
$P_t$ either at a bound ($\pm\bar P$) or determined exactly by an active
$\text{SoC}$ bound, which is the formal reason the verified optimal
schedule is observed to be \emph{bang-bang} (every entry of $P$ equal to
$\pm3$\,kW) except at the one timestep ($P_{10}=-1$\,kW) where the
$\text{SoC}_{\min}=0$ bound binds and pins the power to exactly the
value draining the battery to empty without overshooting -- not an
approximation artifact, but the exact, predicted structure of an LP
vertex solution.
\end{proposition}

\section{Market clearing, congestion pricing, and real options}
\label{sec:merit-order-theory}

\begin{figure}[H]
\centering
\includegraphics[width=0.5\textwidth]{Images/flowchart_market_planning.png}
\caption{Process flow of \pkg{market\_planning.py}.}
\end{figure}

\pkg{market\_\hspace{0pt}planning.py} adds a uniform-price merit-order market
clearing (sorted bids, clearing price read off the marginal bid, not a
solver dual -- see \S\ref{ch:verification}'s dual-sign-bug discussion for
why), a two-node congested-network LMP example, and binomial-lattice
real-option valuation.

\begin{theorem}[Merit-order clearing is LP-optimal, and the clearing price is the dual variable]
\label{thm:merit-order}
For supply bids $(q_i,\pi_i)$ (quantity, price) and demand $d$, the
cost-minimizing dispatch $\min_{0\le x_i\le q_i}\sum_i\pi_ix_i$ s.t.\
$\sum_ix_i=d$ is solved exactly by the merit-order rule (accept bids in
ascending price order until $d$ is met), and the equality constraint's
KKT multiplier $\lambda^\star$ equals the marginal (last-accepted)
bid's price.
\end{theorem}
\begin{proof}
This is a linear program; by Theorem~\ref{thm:strong-duality}, KKT
stationarity for each $x_i$ in the interior of $[0,q_i]$ gives
$\pi_i=\lambda^\star$ (every partially-accepted bid shares the single
market-clearing price $\lambda^\star$), while complementary slackness
forces $x_i=q_i$ (fully accepted) for any bid with $\pi_i<\lambda^\star$
and $x_i=0$ (rejected) for any bid with $\pi_i>\lambda^\star$ -- exactly
the merit-order rule, since any other feasible dispatch could be
strictly improved by an exchange argument (swap one unit from a
higher-priced accepted bid to a lower-priced rejected one, reducing
cost, until no such exchange remains, which is precisely the sorted-
order solution). This is the main course's own Theorem~22, restated here
self-contained for the market module's own two formulations: the plain
merit-order clear (Eq.\ not shown, implemented directly by sorting) and
the congested two-node LMP problem given below, where the per-
node balance constraints $P_a-f=d_a$, $P_b+f=d_b$ each contribute their
own dual variable $\text{LMP}_a,\text{LMP}_b$ by the identical
argument applied node-by-node.
\end{proof}

The two-node LMP problem,
\begin{equation}
\min_{P_a,P_b,f}\; \sum_i \pi_i^a P_a^i + \sum_j \pi_j^b P_b^j
\quad\text{s.t.}\quad P_a-f=d_a,\;\; P_b+f=d_b,\;\; \left|f\right|\le\bar f,
\label{eq:two-node}
\end{equation}
is verified (gen$_a$: 150\,MW@\$18; gen$_b$: 150\,MW@\$48; $d_a=50$,
$d_b=150$, $\bar f=80$) to give $\text{LMP}_a=\$18$,
$\text{LMP}_b=\$48$, spread = $\$30$ (the tie line's shadow price), flow
= $80$\,MW (binding) -- reproduced to the digit in both Python/CVXPY and
MATLAB/\texttt{linprog} \emph{after} fixing a dual-sign bug present in
both (\S\ref{ch:verification}). Note that Theorem~\ref{thm:merit-order}
only identifies the clearing price with \emph{a correctly signed} dual
variable; it says nothing about which raw sign convention a given
solver returns for that dual, which is exactly where both independent
implementations erred (Table~\ref{tab:bugs}).

\begin{theorem}[Real-option value dominates static NPV]
\label{thm:real-option}
Let $V_T$ be the (random) underlying project value at the option's
final exercise date, with payoff $\max(V_T-I,0)$ exercised optimally at
each node of the binomial lattice via the backward recursion given below. Then the
option value $O_{0,0}$ computed by the backward recursion satisfies
$O_{0,0}\ge\max(V_0-I,0)$, the static NPV, with equality iff the
underlying has zero volatility ($u=d=1$).
\end{theorem}
\begin{proof}
At every interior node the recursion takes
$O_{t,j}=\max(\text{exercise now},\text{continuation value})\ge
\text{continuation value}=\frac{p_uO_{t+1,j+1}+(1-p_u)O_{t+1,j}}{1+r}$.
Unrolling this recursion from $t=0$ to $t=T$ along the ``always
continue'' branch gives $O_{0,0}\ge\mathbb
E^{\mathbb Q}\left[\max(V_T-I,0)\right]/(1+r)^T$, the discounted
risk-neutral expectation of exercising only at maturity. Because
$x\mapsto\max(x-I,0)$ is a convex function of $x$, Jensen's inequality
gives $\mathbb E^{\mathbb Q}[\max(V_T-I,0)]\ge\max(\mathbb
E^{\mathbb Q}[V_T]-I,0)=\max(V_0(1+r)^T-I,0)\ge
(1+r)^T\max(V_0-I,0)$ (the last step since $(1+r)^T\ge1$ and the
$\max(\cdot,0)$ is applied to a scaled-up argument), so
$O_{0,0}\ge\max(V_0-I,0)$, the static NPV. Equality throughout requires
Jensen's inequality to be tight, which (for a strictly convex function
of a non-degenerate random variable) happens exactly when $V_T$ is
deterministic, i.e.\ $u=d=1$. This is the main course's own Theorem~23;
restated here because it is the formal justification for the module's
own verified zero-volatility collapse test
(\texttt{test\_\hspace{0pt}binomial\_\hspace{0pt}real\_\hspace{0pt}option\_\hspace{0pt}zero\_\hspace{0pt}volatility\_\hspace{0pt}matches\_\hspace{0pt}static\_\hspace{0pt}npv}),
not merely a convenient coincidence of the chosen test parameters.
\end{proof}

The binomial real-option backward recursion,
\begin{equation}
O_{t,j} = \max\left(V_{t,j}-I,\;\; \frac{p_u O_{t+1,j+1}+(1-p_u)O_{t+1,j}}{1+r}\right),
\qquad V_{t,j}=V_0 u^j d^{t-j},
\label{eq:binomial}
\end{equation}
is verified to give $\text{static\ NPV}=\$10.0$,
$\text{option\ value}=\$27.3590$ on the test parameters $(V_0=100,
I=90, u=1.25, d=0.85, p_u=0.5, n=4, r=0.08)$ -- strictly greater than the
static NPV, confirming Theorem~\ref{thm:real-option}, and confirmed to
collapse exactly to the static NPV at zero volatility ($u=d=1$).

\section{A MOST-equivalent multi-period scheduling tool}

\begin{figure}[H]
\centering
\includegraphics[width=0.5\textwidth]{Images/flowchart_most_scheduling.png}
\caption{Process flow of \pkg{most\_scheduling.py}.}
\end{figure}

\pkg{most\_\hspace{0pt}scheduling.py}/\pkg{psat\_\hspace{0pt}most.m} reproduce the advertised
feature set of MATPOWER's companion MOST tool -- multi-period unit
commitment, ramping, minimum up/down time, storage, reserves, deferrable
demand, N-1 security via LODF, and two-stage stochastic commitment --
which PSAT has no native equivalent of at all. Because neither MILP
solver used (HiGHS via Pyomo; MATLAB's \texttt{intlinprog}) accepts a
true quadratic objective, each generator's convex cost is approximated
by four linear segments of increasing marginal cost; because marginal
cost strictly increases from segment to segment, a cost-minimizing solver
always fills the cheapest segment first with no extra ordering
constraint, a standard exact property of convex piecewise-linear costs.
Verified to within $0.0007\%$ of an exact QP solve.

\begin{proposition}[Piecewise-linear segment filling requires no ordering constraint]
\label{prop:pwl}
Let a generator's cost be approximated by $K$ linear segments with
strictly increasing marginal costs $m_1<m_2<\cdots<m_K$ over consecutive
output ranges, modeled by $K$ non-negative decision variables
$P=\sum_k p_k$, $0\le p_k\le\Delta_k$ (segment width), cost
$\sum_km_kp_k$. Then any cost-minimizing solution automatically
satisfies $p_k=\Delta_k$ (segment $k$ full) for every $k$ with a cheaper
unfilled alternative, i.e.\ segments fill in increasing order of $m_k$,
with no explicit ``fill segment $k$ before $k+1$'' constraint needed.
\end{proposition}
\begin{proof}
Suppose an optimal solution has $p_k<\Delta_k$ (segment $k$ not full)
while $p_{k'}>0$ for some $k'>k$ (a costlier segment partially used).
Since $m_k<m_{k'}$, shifting $\varepsilon=\min(\Delta_k-p_k,p_{k'})>0$
of output from segment $k'$ to segment $k$ keeps $P=\sum_kp_k$ unchanged
(feasible) while changing total cost by $\varepsilon(m_k-m_{k'})<0$ -- a
strict improvement, contradicting optimality. Hence no such pair
$(k,k')$ can exist at an optimum, which is exactly the claimed ordered-
filling property. This is a direct corollary of the problem's convexity
(strictly increasing marginal cost is precisely the condition for a
piecewise-linear cost function to be convex) and requires no integer
ordering variables, the reason \pkg{most\_\hspace{0pt}scheduling.py}'s MILP
formulation needs binary variables only for commitment, never for
segment order.
\end{proof}

\textbf{Cross-language validation}: on a hand-traceable cycling-economics
scenario (a baseload unit plus a pricier unit with a \$200 startup cost
and \$5/period no-load cost, demand oscillating $[40,\allowbreak70,\allowbreak40,\allowbreak70,\allowbreak40]$\,MW),
both the Python/HiGHS and MATLAB/\texttt{intlinprog} solvers -- built
completely independently -- reproduce the identical commitment pattern
$u=[0,1,1,1,0]$ and total cost \$3615.0000 to the penny.

\textbf{Security constraint, with a genuine LODF bug found and fixed.}
The Line Outage Distribution Factor,
\begin{equation}
\text{LODF}_{l,k} = \frac{\text{PTDF}_{l,i}-\text{PTDF}_{l,j}}
{1-(\text{PTDF}_{k,i}-\text{PTDF}_{k,j})},
\label{eq:lodf}
\end{equation}
had an early draft with a spurious extra reactance factor in the
numerator/denominator, silently under-constraining the dispatch: on a
3-bus system with a stated 40\,MW post-contingency limit, the resulting
dispatch actually permitted 58\,MW, caught only by independently
re-solving the post-contingency network under the chosen dispatch. Fixed
to the standard textbook formula given above; the identical
re-solve afterward gives exactly 40.0\,MW.

\begin{theorem}[LODF from the Sherman--Morrison rank-one outage update]
\label{thm:lodf-derivation}
Outaging line $k=(p,q)$ with pre-outage admittance $y_k=1/x_k$ is
equivalent, in the DC model, to adding a rank-one correction
$-y_k(e_p-e_q)(e_p-e_q)^\top$ to the Laplacian $B'$ (removing exactly
that line's contribution). The post-outage reduced Laplacian is
therefore $B'_{\text{post}}=B'-y_k ww^\top$ with $w=e_p-e_q$, and
Theorem~\ref{thm:sherman-morrison} gives
\begin{equation}
(B'_{\text{post}})^{-1} = (B')^{-1} + \frac{y_k(B')^{-1}ww^\top(B')^{-1}}{1-y_kw^\top(B')^{-1}w}.
\label{eq:lodf-sm}
\end{equation}
Combined with $\text{PTDF}=\text{diag}(y)A(B')^{-1}$ (branch flow as
a linear function of injections, $A$ the reduced incidence matrix), the
change in any monitored line $l$'s flow caused by outaging $k$, divided
by line $k$'s pre-outage flow, reduces exactly to the standard LODF
formula stated earlier
-- with \emph{no} extra factor of $x_k$ anywhere in the final expression,
because the $y_k$ appearing in the rank-one update numerator and the
$y_k$ implicit in $\text{PTDF}_{k,\cdot}$'s own definition cancel
algebraically.
\end{theorem}
\begin{proof}[Proof sketch]
Two substitutions turn the rank-one update above into the LODF formula
stated earlier. First,
by definition $\text{PTDF}_{k,i}-\text{PTDF}_{k,j}=y_k\,e_k^\top
A(B')^{-1}w=y_k\,w^\top(B')^{-1}w$ (using $A^\top e_k=w$ for the branch-$k$
row of the incidence matrix, and $(B')^{-1}$ symmetric), so the
denominator of the rank-one update, $1-y_kw^\top(B')^{-1}w$, is exactly
$1-(\text{PTDF}_{k,i}-\text{PTDF}_{k,j})$ -- matching
the LODF formula's own denominator with no leftover factor of $x_k$ or
$y_k$, since the single $y_k$ introduced by the Sherman--Morrison
numerator is exactly the $y_k$ already implicit in the definition of
$\text{PTDF}_{k,i}-\text{PTDF}_{k,j}$ itself. Second, propagating
the same substitution through $\text{PTDF}_{l,\cdot}
(B'_{\text{post}})^{-1}$'s numerator correction term and dividing by
line $k$'s own pre-outage flow (which carries the \emph{same} prefactor
$y_k\,e_l^\top A(B')^{-1}w$ structure for the monitored line $l$)
cancels every remaining $y_k$, leaving exactly the ratio of
$\text{PTDF}$ differences the LODF formula states. The complete
line-by-line bookkeeping is standard; see Wood, Wollenberg \&
Sheblé~\cite{woodwollenberg}, \S11.2, for the full derivation. This
confirms, by direct algebraic derivation -- not only by the independent
post-contingency re-solve this project actually used to catch the bug
-- that the standard LODF formula is the unique algebraically correct formula,
and that the spurious extra reactance factor in the first draft could
not have arisen from a correct application of the Sherman--Morrison
identity.
\end{proof}

\section{Stability-constrained optimal power flow}
\label{sec:stability-opf-theory}

\begin{figure}[H]
\centering
\includegraphics[width=0.5\textwidth]{Images/flowchart_stability_constrained_opf.png}
\caption{Process flow of \pkg{stability\_constrained\_opf.py}.}
\end{figure}

\pkg{stability\_\hspace{0pt}constrained\_\hspace{0pt}opf.py}/\pkg{psat\_\hspace{0pt}stability\_\hspace{0pt}constrained\_\hspace{0pt}
opf.m} couple a DC-OPF dispatch to a genuine multi-machine small-signal
stability check -- none of PSAT's dispatch-side tools touch dynamic
eigenvalues at all. Linearizing the swing equation about an operating
point $\theta^\star$ gives
\begin{equation}
A = \begin{bmatrix}0 & I \\ -M^{-1}K(\theta^\star) & -M^{-1}D\end{bmatrix},
\qquad K_{ij}(\theta^\star)=B_{ij}\cos(\theta_i^\star-\theta_j^\star),
\label{eq:stability-A}
\end{equation}
with $M=\text{diag}(2H_i)$. Because $\cos(\cdot)$ changes sign exactly
at a $90^\circ$ angle difference, the system's stability undergoes a
genuine discontinuous transition there, not a gradual degradation --
verified directly: sweeping one line's angle difference through
$89^\circ\to90^\circ\to91^\circ$ gives maximum eigenvalue real part
$-0.0500\to-0.0500\to+0.0466$ (stable, boundary, unstable). For a
uniformly $H$/$D$-parameterized system, the eigenvalue real part stays
pinned at exactly $-D/(4H)$ for every stable configuration, regardless of
$K$'s magnitude, until the $90^\circ$ crossing itself.

\begin{theorem}[Multi-machine linearized stability, generalizing Theorem~\ref{thm:vsm}]
\label{thm:multimachine-stability}
$K(\theta^\star)\in\mathbb{R}^{n\times n}$ in the linearized state
matrix above is
itself a weighted graph Laplacian (Definition in
\S\ref{sec:graph-prelim}) with edge weights $B_{ij}\cos(\theta_i^\star-
\theta_j^\star)$, which are positive exactly for lines with angle
difference under $90^\circ$. If every line's angle difference satisfies
$\left|\theta_i^\star-\theta_j^\star\right|<90^\circ$ (every edge weight positive),
then by Theorem~\ref{thm:laplacian}, $K(\theta^\star)$ is PSD with a
single zero eigenvalue (eigenvector $\mathbf1$, the physically arbitrary
common angle reference), and the linearized system $A$ given
above has exactly one trivial zero eigenvalue (from the
reference-angle mode) and every other eigenvalue in the open left half-
plane whenever $D_i>0$ for every $i$. If instead some line's angle
difference exceeds $90^\circ$, $K(\theta^\star)$ generically loses
positive semi-definiteness (it becomes an \emph{indefinite}, signed
Laplacian) and $A$ generically acquires at least one eigenvalue with
strictly positive real part.
\end{theorem}
\begin{proof}
Write the swing-equation energy function $W(\delta,\omega)=
\tfrac12\omega^\top M\omega+U(\delta)$ with potential
$U(\delta)=-\sum_{i<j}B_{ij}\cos(\delta_i-\delta_j)$ (the standard
power-system transient-energy function, Kundur~\cite{kundur}, Ch.~12).
Its Hessian in $\delta$ at $\delta^\star$ is exactly $K(\theta^\star)$:
$\partial^2U/\partial\delta_i\partial\delta_j=-B_{ij}\cos(\theta_i^\star-
\theta_j^\star)$ off-diagonal (matching $K$'s off-diagonal entries up to
the sign convention absorbed into $A$) and
$\partial^2U/\partial\delta_i^2=\sum_{j}B_{ij}\cos(\theta_i^\star-
\theta_j^\star)$ on the diagonal, which is exactly a weighted Laplacian
with weights $B_{ij}\cos(\theta_i^\star-\theta_j^\star)$. When every
weight is positive, Theorem~\ref{thm:laplacian} applies directly (the
proof there uses only positivity of the edge weights, not that they
equal a physical susceptance): $K(\theta^\star)\succeq0$ with the
single null direction $\mathbf1$, meaning $\theta^\star$ is a local
minimum of $U$ along every direction except the physically meaningless
uniform angle shift -- the standard Lyapunov argument for swing-equation
stability (energy function has a local minimum at the equilibrium) then
gives local asymptotic stability of the full $(\delta,\omega)$ system
once positive damping is added, matching the single-machine case's
Theorem~\ref{thm:vsm} exactly but now for $n$ coupled machines. Once a
weight turns negative, $K(\theta^\star)$ is no longer guaranteed
PSD and $\theta^\star$ is generically a saddle point of $U$, not a
local minimum, matching the observed sign flip in the eigenvalue
sweep.
\end{proof}

\begin{corollary}[Angle-difference limits are a conservative convex certificate for stability]
\label{cor:angle-certificate}
A per-line angle-difference limit of $40$--$70^\circ$ (strictly inside
the $90^\circ$ threshold, as used throughout
\pkg{stability\_\hspace{0pt}constrained\_\hspace{0pt}opf.py}) is \emph{sufficient}
for the stability condition of Theorem~\ref{thm:multimachine-stability}
(every weight stays positive, with margin), but not formally
\emph{necessary} for stability in general (a system could in principle
have one line beyond $90^\circ$ while the aggregate Laplacian structure
still happens to stay PSD on the relevant subspace). This asymmetry --
a tractable convex sufficient condition standing in for an intractable
exact one -- is the same structural trade-off every convex
restriction of a non-convex feasible region makes, and is stated here
explicitly rather than left implicit, consistent with
\S\ref{ch:limitations}'s honesty discipline.
\end{corollary}

A per-line angle-difference limit is used as a convex, OPF-embeddable
proxy for stability (angle-difference bounds are linear in the decision
variables, unlike a direct eigenvalue-sign constraint, which is not a
convex -- or even explicit -- function of dispatch), with an iterative
cutting-plane loop adding a binding constraint for any violating line
until none remain. On a 4-bus ring test system (branch reactances
$[1.8,0.3,0.3,0.3]$\,pu, demand $[0,140,120,0]$\,MW): the unconstrained
economic dispatch concentrates generation at the cheap slack bus
($P_g=[260,0,0,0]$, cost \$2600.00, worst angle $75.63^\circ$, already
stable); a $50^\circ$ limit forces $P_g=[173.17,50.0,36.83,0]$, cost
\$5205.02 -- a \textbf{price of stability of +100.19\%} -- verified
monotonically non-decreasing as the limit tightens, and verified to
report infeasibility honestly (rather than crash) when no feasible
dispatch exists under a sufficiently tight limit.

%=============================================================
\chapter{Security and Resilience Extensions}
\label{ch:resilience}
%=============================================================

\section{DC state estimation and stealthy false-data injection}
\label{sec:fdia-theory}

\begin{figure}[H]
\centering
\includegraphics[width=0.5\textwidth]{Images/flowchart_fdia_state_estimation.png}
\caption{Process flow of \pkg{fdia\_state\_estimation.py}.}
\end{figure}

\pkg{fdia\_\hspace{0pt}state\_\hspace{0pt}estimation.py} builds a DC state-estimation measurement
model (bus-injection and branch-flow measurements) on PYPOWER's real
\texttt{case9} network, extending PSAT's own PMU-placement
(\texttt{fm\_pmuloc}) and state-estimation (\texttt{fm\_lssest}) routines
with the bad-data-detection and attack-construction capability neither
has. The weighted-least-squares estimate and chi-square detection
statistic are
\begin{equation}
\hat x = (H^\top R^{-1}H)^{-1}H^\top R^{-1}z, \qquad
J = r^\top R^{-1}r,\quad r=z-H\hat x,
\label{eq:se}
\end{equation}
flagged as bad data when $J$ exceeds the $\chi^2_{1-\alpha}(m-n)$
threshold -- the exact test Theorem~\ref{thm:chisquare} derives,
specialized to $k=m-n$ degrees of freedom for this WLS residual.

\begin{theorem}[Exact characterization of undetectable (stealthy) attacks]
\label{thm:fdia-exact}
Let $z_a=z+a$ be the attacked measurement vector for some attack vector
$a\ne0$, giving attacked estimate $\hat x_a=(H^\top R^{-1}H)^{-1}H^\top
R^{-1}z_a$ and residual $r_a=z_a-H\hat x_a$. Then $r_a=r$ (the residual,
and hence $J$, is \emph{exactly} unchanged by the attack) \emph{if and
only if} $a\in\text{range}(H)$, i.e.\ $a=Hc$ for some $c\in
\mathbb{R}^n$ -- and in that case the attacked state estimate shifts by
exactly $\hat x_a=\hat x+c$, the attacker's chosen vector.
\end{theorem}
\begin{proof}
($\Leftarrow$) If $a=Hc$: $\hat x_a=(H^\top R^{-1}H)^{-1}H^\top
R^{-1}(z+Hc)=\hat x+(H^\top R^{-1}H)^{-1}H^\top R^{-1}Hc=\hat x+c$, so
$r_a=z+Hc-H(\hat x+c)=z-H\hat x=r$ exactly, by direct cancellation --
the computation given informally in the main text, now exhibited as a
genuine ``if'' direction. ($\Rightarrow$) Conversely, suppose $r_a=r$.
Define the residual (hat) matrix $S=I-H(H^\top R^{-1}H)^{-1}H^\top
R^{-1}$, an oblique projector satisfying $r=Sz$ for every $z$ (idempotent,
$S^2=S$, with $\text{range}(H)=\ker(S)$ -- direct verification from
the definition of $S$). Then $r_a=S(z+a)=r+Sa$, so $r_a=r$ forces
$Sa=0$, i.e.\ $a\in\ker(S)=\text{range}(H)$, establishing $a=Hc$ for
some $c$. Thus $a=Hc$ is not merely \emph{a} family of stealthy attacks
but \emph{exactly} the set of perfectly undetectable ones -- this is the
main course's own Theorem~19, restated here with the full both-
directions proof (the main text's informal derivation only showed the
easy direction); see also Liu, Ning \&
Reiter~\cite{liu}, \S4, for the original construction.
\end{proof}

Reusing this course's own already-verified case9 setup, the stealthy
attack leaves $J$ identical to the clean case to the numerical noise
floor while shifting the state estimate by exactly the attacker's chosen
vector, confirmed robust across 10 independent random seeds, whereas a
random (non-stealthy) bad-data injection of comparable magnitude is
correctly flagged by the same test -- exactly the dichotomy
Theorem~\ref{thm:fdia-exact} predicts, since a generic random
perturbation lies in $\text{range}(H)$ only with probability zero
(a measure-zero $n$-dimensional subspace of the $m$-dimensional
measurement space, $m>n$).

\section{Extended attack surface: GPS spoofing and a measurement-protection game}

\begin{figure}[H]
\centering
\includegraphics[width=0.5\textwidth]{Images/flowchart_attack_surface.png}
\caption{Process flow of \pkg{attack\_surface.py}.}
\end{figure}

\pkg{attack\_\hspace{0pt}surface.py} adds two further, mechanism-distinct attack
dimensions. \textbf{GPS spoofing} on a PMU's time reference shifts its
reported phase angle by
\begin{equation}
\Delta\theta = 2\pi f_{\text{nominal}}\,\Delta t,
\label{eq:gps-spoof}
\end{equation}
for a clock offset $\Delta t$; a single spoofed PMU generically fails to
lie in $\text{range}(H)$ and is (correctly) flagged by the same
chi-square test, but an attacker able to spoof enough clocks
simultaneously can choose each $\Delta t$ to realize \emph{any} $a=Hc$
attack exactly -- GPS spoofing is shown to be an alternative physical
delivery mechanism for the identical stealthy construction, not a new
detection loophole.

The \textbf{measurement-protection game} formalizes an empirical finding
from this course's own teaching content (protecting 8 of 17 measurements
blocks every stealthy attack) as the exact linear-algebra condition it
is.

\begin{theorem}[Measurement-protection rank condition]
\label{thm:protection-rank}
Partition the measurements into a protected set $P$ (rows the attacker
cannot alter) and an unprotected set $U$, with $H=\begin{bmatrix}H_P\\
H_U\end{bmatrix}$ after reordering. An attacker confined to $U$ (i.e.\
$a_P=0$) can realize a nonzero stealthy attack ($a=Hc\ne0$ with
$a_P=0$) if and only if $H_P$ does \emph{not} have full column rank $n$.
Equivalently, protecting $P$ blocks \emph{every} possible stealthy
attack if and only if $\text{rank}(H_P)=n$.
\end{theorem}
\begin{proof}
($\Leftarrow$, contrapositive) If $\text{rank}(H_P)<n$, then $H_P$ has
a nontrivial null space: there exists $c\ne0$ with $H_Pc=0$. Taking
$a=Hc$ (a valid stealthy attack by Theorem~\ref{thm:fdia-exact}) gives
$a_P=H_Pc=0$ exactly -- the attacker needs to alter no protected
measurement at all, realizing a nonzero attack using only unprotected
measurements in $U$ (since $a_U=H_Uc$ need not vanish for a generic
$c$ in $\ker(H_P)$ outside $\ker(H)$, and $\ker(H)=\{0\}$ whenever $H$
itself has full column rank, the standing observability assumption).
($\Rightarrow$, contrapositive) If $\text{rank}(H_P)=n$, then
$H_Pc=0\Rightarrow c=0$ (full column rank means trivial null space), so
$a_P=H_Pc=0$ forces $c=0$ and hence $a=Hc=0$ -- no nonzero stealthy
attack can avoid touching $P$. Both directions together give the
claimed ``if and only if''. See also the parallel treatment in
Kosut et al.\ and the main course's own Theorem~19's corollary on
measurement protection.
\end{proof}

\begin{corollary}[Minimum protection size is governed by $H$'s column structure]
\label{cor:min-protection}
The smallest protection set $P$ satisfying Theorem~\ref{thm:protection-rank}
has $\left|P\right|\ge n$ (a rank-$n$ matrix needs at least $n$ rows), with
equality achievable if and only if some $n$ rows of $H$ are linearly
independent and span a complement making the remaining structure
observable -- generically true for a well-conditioned $H$, but the
\emph{specific} rows that work depend on $H$'s own structure, which is
exactly why this module's greedy heuristic (adding, at each step, the
unprotected measurement whose addition most increases $\text{rank}
(H_P)$) is not guaranteed optimal: greedy rank-augmentation is a
matroid-type problem for which the greedy algorithm IS optimal under a
matroid independence structure, but the \emph{specific} objective here
(minimum cardinality spanning set) coincides with matroid-greedy
optimality only when no additional combinatorial cost differentiates
candidate rows -- confirmed directly by this module's own brute-force
comparison finding cases where greedy is suboptimal.
\end{corollary}

A greedy defense heuristic is built and directly compared against
brute-force optimal, confirming the 8-measurement minimum defense and
the under-10-measurement zero-attacker-freedom threshold predicted by
Theorem~\ref{thm:protection-rank} and Corollary~\ref{cor:min-protection},
and showing greedy is not always optimal.

\section{Cascading failure and N-k analysis}

\begin{figure}[H]
\centering
\includegraphics[width=0.45\textwidth]{Images/flowchart_cascading_failure.png}
\caption{Process flow of \pkg{cascading\_failure.py}.}
\end{figure}

\pkg{cascading\_\hspace{0pt}failure.py} extends the single-level N-1 screen of
Chapter~\ref{ch:architecture} to a sequential cascade: after every line
trip, the DC power-flow model $P=B\theta$ is re-solved, and the single
most-overloaded remaining in-service line ($\left|\text{flow}\right|>\text{factor}\times\text{rate}A$,
where $\text{factor}>1$ is the overload-trigger multiplier) trips next, stopping when no line is
overloaded or the network islands. Verified on a specially constructed
``overload-rich'' test system to produce a genuine, hand-checkable 2-step
cascade; an honest null result is also reported -- PYPOWER's
\texttt{case118} did \emph{not} cascade under the overload factors
tested, reported as found rather than adjusted to manufacture a more
dramatic result.

\section{Decentralized wide-area frequency control}

\begin{figure}[H]
\centering
\includegraphics[width=0.46\textwidth]{Images/flowchart_wide_area_control.png}
\caption{Process flow of \pkg{wide\_area\_control.py}.}
\end{figure}

\pkg{wide\_\hspace{0pt}area\_\hspace{0pt}control.py} generalizes PSAT's own single/two-area AGC
content to $K>2$ areas with no central coordinator, each area's own AGC
integrator driven by a Metropolis--Hastings weighted consensus average of
its neighbors' area control error,
\begin{equation}
\text{ACE}_k = \left(P_{\text{tie},k}^{\text{actual}}-P_{\text{tie},k}^{\text{sched}}\right)
+ \beta_k\,(f-f_{\text{nominal}}).
\label{eq:ace}
\end{equation}
\textbf{A genuine instability was found and fixed}: applying the
consensus weight matrix only once per physical timestep is unstable, not
merely slow -- a 4-area case's frequency deviation stayed under
0.01\,Hz through $t=150$\,s, then diverged to $29{,}821.0$\,Hz by
$t=600$\,s, because the AGC integrator has no mechanism to ``forget'' the
small per-step averaging error a single consensus application leaves
behind, which compounds through the integrator over time. Fixed by
applying the consensus matrix repeatedly (20 sub-iterations) per physical
step, after which the maximum deviation at $t=300$\,s is $1.06\times
10^{-5}$\,Hz, matching the centralized baseline to that same tolerance.

%=============================================================
\chapter{Control Extensions}
\label{ch:control}
%=============================================================

\section{MPC and LQR from a PSAT linearization}
\label{sec:mpc-theory}

\begin{figure}[H]
\centering
\includegraphics[width=0.46\textwidth]{Images/flowchart_mpc_export.png}
\caption{Process flow of \pkg{mpc\_export.py}.}
\end{figure}

\pkg{mpc\_\hspace{0pt}export.py}/\pkg{psat\_\hspace{0pt}discretize.m}+\pkg{psat\_\hspace{0pt}mpc\_\hspace{0pt}design.m}
turn Chapter~\ref{ch:architecture}'s $(A,B,C,D)$ export into a finite-
horizon MPC and an infinite-horizon LQR baseline.

\begin{theorem}[Exact zero-order-hold discretization]
\label{thm:zoh}
For $\dot x=Ax+Bu$ with $u(t)\equiv u_k$ held constant for
$t_k\le t<t_k+\Delta t$, the exact state at $t_k+\Delta t$ is
$x(t_k+\Delta t)=A_dx(t_k)+B_du_k$ with $A_d=e^{A\Delta t}$,
$B_d=\left(\int_0^{\Delta t}e^{A\tau}d\tau\right)B$, and
the augmented-matrix-exponential formula above computes both simultaneously via a single matrix
exponential of the augmented block matrix.
\end{theorem}
\begin{proof}
The variation-of-constants formula gives $x(t_k+\Delta t)=
e^{A\Delta t}x(t_k)+\int_0^{\Delta t}e^{A(\Delta t-\tau)}Bu_k\,d\tau=
e^{A\Delta t}x(t_k)+\left(\int_0^{\Delta t}e^{As}ds\right)Bu_k$ (substituting
$s=\Delta t-\tau$), matching $A_d,B_d$ above. For the augmented matrix
exponential, direct block computation gives
$\exp\left(\begin{bmatrix}A&B\\0&0\end{bmatrix}\Delta t\right)=
\begin{bmatrix}e^{A\Delta t} & \int_0^{\Delta t}e^{As}ds\,B\\ 0 &
I\end{bmatrix}$ (the $(2,2)$ block exponentiates the zero matrix to
$I$; the $(1,2)$ block is exactly the integral series for $B_d$, since
the augmented matrix's own power series reproduces the Duhamel/
variation-of-constants integral term-by-term) -- see Van Loan's classic
1978 note on computing matrix exponentials for state-space
discretization, or Franklin, Powell \& Workman, \emph{Digital Control
of Dynamic Systems}, \S4.3. This is exact for any $\Delta t$ (no
truncation error), unlike forward-Euler's $A_d\approx I+A\Delta t$,
which is only first-order accurate.
\end{proof}

The finite-horizon MPC solves a quadratic program in $(x,u)$ with
weights $Q,R_w$ (a direct application of Theorem~\ref{thm:riccati}'s
backward recursion, truncated to a finite horizon with terminal weight
$P_N$); the infinite-horizon LQR baseline is obtained from the discrete
algebraic Riccati equation (DARE), the $N\to\infty$ fixed point of that
same recursion.

\begin{theorem}[No-preview finite-horizon MPC converges to infinite-horizon LQR]
\label{thm:mpc-lqr}
For a disturbance-free, unconstrained finite-horizon LQR problem with
terminal weight $P_N=Q$ (or any $P_N\succeq0$), the Riccati recursion
given earlier run backward from $P_N$ converges, as the
horizon $N\to\infty$, to the unique stabilizing solution $P_\infty$ of
the DARE, and the resulting first-step gain $K_0\to K_{\text{lqr}}=
(R_w+B_d^\top P_\infty B_d)^{-1}B_d^\top P_\infty A_d$.
\end{theorem}
\begin{proof}
The Riccati recursion given earlier, run backward, is
exactly the Riccati \emph{difference} equation iterated forward in the
reversed index $j=N-k$; under stabilizability of $(A_d,B_d)$ and
detectability of $(A_d,Q^{1/2})$ (Theorem~\ref{thm:riccati}'s standing
hypotheses), this iteration is a monotone, bounded sequence of PSD
matrices (monotone because adding one more stage to the horizon can only
weakly increase the achievable cost-to-go's curvature; bounded above
because the DARE's stabilizing solution provides a fixed upper bound,
and below by $0$), hence convergent by the monotone convergence theorem
for matrix sequences in the Loewner order; the limit must be a fixed
point of that same recursion, i.e.\ a solution of the DARE, and
standard DARE theory (Anderson \& Moore~\cite{andersonmoore}, Ch.~2)
identifies the limit of this particular (monotone, from $P_N\succeq0$)
sequence as the \emph{stabilizing} solution specifically, not one of the
DARE's other (non-stabilizing) solutions. This formally explains the
module's own verified numeric convergence
($N=1$: $K_0=-1.319044$; $N\ge5$: converged to $-1.410434=
-K_{\text{lqr}}$ computed independently via
\texttt{scipy.linalg.solve\_\hspace{0pt}discrete\_\hspace{0pt}are}) as an instance of a
general theorem, not a property specific to the one test system used.
\end{proof}

Verified: a no-disturbance-preview finite-horizon MPC's first-step
control converges smoothly as the horizon $N$ grows ($N=1$: $-1.319044$;
$N\ge5$: $-1.410434$, stable) to exactly $-K_{\text{lqr}}x_0=-1.410434$,
where $K_{\text{lqr}}=1.41043436$ is computed independently from the
DARE -- a genuine cross-method numeric match isolating disturbance
preview, not horizon length per se, as the real source of MPC's advantage
over plain LQR in this course's own worked example (the comparison is
fair only because Theorem~\ref{thm:mpc-lqr} guarantees both methods
converge to the \emph{identical} control law in the no-disturbance
limit, isolating disturbance preview as the one remaining difference).

\S\ref{ch:resilience}'s wide-area control and this chapter's MPC/LQR
work share one methodological thread: both verify a control design by
checking convergence to an \emph{independently computed} fixed point
(the centralized baseline; the DARE gain) rather than only checking that
the closed loop is numerically stable.

%=============================================================
\chapter{Learning-Based Methods}
\label{ch:learning}
%=============================================================

\section{A graph neural network surrogate for AC power flow}

\begin{figure}[H]
\centering
\includegraphics[width=0.5\textwidth]{Images/flowchart_gnn_power_flow.png}
\caption{Process flow of \pkg{gnn\_power\_flow.py}.}
\end{figure}

\pkg{gnn\_\hspace{0pt}power\_\hspace{0pt}flow.py} trains a plain-PyTorch (no
\texttt{torch\_geometric}, unavailable in this environment) message-
passing network on real PYPOWER Newton--Raphson AC solutions. One
message-passing layer computes, for every directed edge $u\to v$,
\begin{equation}
\text{msg}_{uv} = f_\theta\left([x_u,x_v,e_{uv}]\right), \qquad
\text{agg}_v = \sum_{u:(u,v)\in E}\text{msg}_{uv}, \qquad
x_v' = g_\phi\left([x_v,\text{agg}_v]\right),
\label{eq:gnn}
\end{equation}
implemented as a batched \texttt{index\_add\_} scatter-sum, functionally
equivalent to a basic \texttt{torch\_geometric} \texttt{MessagePassing}
layer with sum aggregation.

\begin{proposition}[Permutation equivariance of message passing]
\label{prop:gnn-equivariance}
Let $\pi$ be any permutation of the node index set. If the node/edge
features are relabeled consistently ($x_{\pi(v)}\leftarrow x_v$,
$e_{\pi(u)\pi(v)}\leftarrow e_{uv}$), the message-passing layer
given earlier produces outputs related by the same permutation,
$x'_{\pi(v)}(\text{relabeled input})=x'_v(\text{original input})$.
\end{proposition}
\begin{proof}
The aggregation $\text{agg}_v=\sum_{u:(u,v)\in E}\text{msg}_{uv}$ is a
sum over the (unordered) in-neighborhood of $v$, and summation is
commutative/associative, so it does not depend on any particular
ordering of the neighbors -- relabeling nodes by $\pi$ merely relabels
which summands belong to which output index, without changing any
individual message's value (since $f_\theta$ is applied with the same
weights to every edge, regardless of node identity). The update
$x_v'=g_\phi([x_v,\text{agg}_v])$ then inherits the same equivariance.
This is the standard graph-neural-network permutation-equivariance
property (see, e.g., Bronstein et al., \emph{Geometric Deep Learning},
\S2.2--4, for the general theory); it is the formal reason a single
trained model can be applied unchanged to any bus numbering/ordering of
a \emph{fixed} topology (what makes batched training over many randomly
sampled operating points of \texttt{case14} meaningful at all), while
offering no guarantee whatsoever across \emph{different} topologies
(different edge sets $E$, different node counts) -- exactly the
boundary this module's own honest cross-topology negative result
(below) falls on.
\end{proof}

\textbf{A genuine training bug was found and
fixed}: an un-normalized shared loss over both outputs (voltage magnitude
and angle) let angle error -- with roughly $14\times$ the squared-error
magnitude spread of voltage magnitude -- dominate every gradient step,
producing a model that scored \emph{worse} on voltage-magnitude accuracy
(0.0044\,pu MAE) than a trivial frozen-base-case baseline (0.0014\,pu),
even though it clearly beat that baseline on angle. Fixed by z-score
normalizing each output before computing the loss: the identical
architecture then beats the frozen baseline on \emph{both} outputs
(0.00137 vs.\ 0.00144\,pu on magnitude; 0.00876 vs.\ 0.01470\,rad on
angle), with roughly $270\times$ faster batched inference than re-solving
via Newton--Raphson. An honest negative result is also reported: a model
trained on \texttt{case14} and evaluated zero-shot on \texttt{case30}
performs far worse than on its own held-out operating points --
surrogate models of this kind are topology-specific, not transferable, a
genuinely useful pedagogical finding rather than a glossed-over
limitation.

\section{A reinforcement-learning battery-scheduling environment}
\label{sec:rl-theory}

\begin{figure}[H]
\centering
\includegraphics[width=0.46\textwidth]{Images/flowchart_rl_env.png}
\caption{Process flow of \pkg{rl\_env.py}.}
\end{figure}

\pkg{rl\_\hspace{0pt}env.py} provides a Gymnasium-compatible battery-scheduling
environment with an independent exact dynamic-programming optimum as
ground truth (cross-checked against full brute-force enumeration), the
same style of honest benchmarking Chapter~\ref{ch:optimization} applies
to every optimization-under-uncertainty tool in this report.

\begin{theorem}[Bellman optimality and the contraction property of the Bellman operator]
\label{thm:bellman}
For a Markov decision process with discount $0\le\gamma_{\text{RL}}<1$,
bounded reward $r$, the optimal value function $V^\star$ satisfies the
Bellman optimality equation
$V^\star(s)=\max_a\left[r(s,a)+\gamma_{\text{RL}}\mathbb E_{s'}[V^\star(s')]\right]$,
and the Bellman operator $(\mathcal TV)(s)\triangleq\max_a\left[r(s,a)+
\gamma_{\text{RL}}\mathbb E_{s'}[V(s')]\right]$ is a contraction in the
sup-norm with modulus $\gamma_{\text{RL}}$, so by
Theorem~\ref{thm:banach} it has a unique fixed point (necessarily
$V^\star$) and value iteration $V_{k+1}=\mathcal TV_k$ converges to it
geometrically from any initialization.
\end{theorem}
\begin{proof}
For any two value functions $V_1,V_2$ and any state $s$, let
$a_1^\star=\arg\max_aQ_1(s,a)$ where $Q_i(s,a)=r(s,a)+
\gamma_{\text{RL}}\mathbb E_{s'}[V_i(s')]$. Then $(\mathcal
TV_1)(s)-(\mathcal TV_2)(s)\le Q_1(s,a_1^\star)-Q_2(s,a_1^\star)=
\gamma_{\text{RL}}\mathbb E_{s'}[V_1(s')-V_2(s')]\le
\gamma_{\text{RL}}\|V_1-V_2\|_\infty$ (bounding the expectation by
the sup-norm difference); by symmetry $\left|(\mathcal TV_1)(s)-(\mathcal
TV_2)(s)\right|\le\gamma_{\text{RL}}\|V_1-V_2\|_\infty$ for every $s$,
giving $\|\mathcal TV_1-\mathcal TV_2\|_\infty\le\gamma_{\text{RL}}
\|V_1-V_2\|_\infty$, the contraction property with modulus
$\gamma_{\text{RL}}<1$. The space of bounded functions on the (finite,
here) state space with the sup-norm is complete, so
Theorem~\ref{thm:banach} applies directly, giving both existence/
uniqueness of $V^\star$ and geometric convergence of value iteration.
This is the main course's own Theorem~18 (there attributed to Watkins \&
Dayan's Q-learning convergence proof, which applies the identical
contraction argument to the state-action value function $Q$ rather than
$V$); restated here because it is the formal reason this module's own
tabular Q-learning benchmark provably converges to the exact DP optimum
(reward = $60.0$ on this project's reference scenario) given enough
exploration, independent of the specific exploration schedule used, as
long as every state-action pair is visited infinitely often (the
standard stochastic-approximation condition Watkins \& Dayan add on top
of the deterministic contraction argument above).
\end{proof}

It reproduces this course's own documented reward-hacking finding:
removing the terminal-state-of-charge bonus term from the reward
measurably drains the battery rather than landing it on the intended
target -- a clean illustration that Theorem~\ref{thm:bellman} guarantees
convergence to the optimum \emph{of whatever reward function is
specified}, not to any behavior the designer actually intended; the
theorem is silent on reward design, which is exactly where this
particular failure originates (objective mis-specification, not an
algorithmic failure).

%=============================================================
\chapter{Uncertainty Quantification}
\label{ch:uncertainty}
%=============================================================

\section{Probabilistic power flow}
\label{sec:ppf-theory}

\begin{figure}[H]
\centering
\includegraphics[width=0.5\textwidth]{Images/flowchart_probabilistic_power_flow.png}
\caption{Process flow of \pkg{probabilistic\_power\_flow.py}.}
\end{figure}

\pkg{probabilistic\_\hspace{0pt}power\_\hspace{0pt}flow.py}/\pkg{psat\_\hspace{0pt}probabilistic\_\hspace{0pt}power\_\hspace{0pt}
flow.m} quantify uncertainty \emph{in a power-flow solution itself}
(given uncertain bus injections, what is the distribution of line flows?)
-- a distinct question from Chapter~\ref{ch:optimization}'s
robust/stochastic \emph{dispatch} tools, which decide what to dispatch
under uncertainty.

\begin{corollary}[Exactness of analytical PPF for DC power flow]
\label{cor:ppf-exact}
If $\text{injections}\sim\mathcal N(\mu_{\text{inj}},\Sigma_{\mathrm
{inj}})$, then since $\text{flow}=\text{PTDF}\cdot\text{injections}$ is
an affine (indeed linear) function of a Gaussian vector,
Theorem~\ref{thm:affine-gaussian} gives \emph{exactly}
$\text{flow}\sim\mathcal N(\text{PTDF}\,\mu_{\text{inj}},\;
\text{PTDF}\,\Sigma_{\text{inj}}\,\text{PTDF}^\top)$ -- the analytical
covariance formula given below
is not an approximation, moment-matching heuristic, or linearization: it
is the exact distribution, with no error term, for \emph{any} DC
network and \emph{any} Gaussian injection distribution. This is why the
analytical method in this module is validated by showing it is the
limit Monte Carlo \emph{converges toward} (\S\ref{ch:uncertainty}'s own
Task~1-style verification), rather than the usual relationship between
an approximation and a ground truth.
\end{corollary}
\begin{proof}
Direct application of Theorem~\ref{thm:affine-gaussian} with
$A=\text{PTDF}$, $b=0$; no further argument is needed because DC power
flow's linearity (not merely approximate linearity, since the DC model
itself is already a linearization of AC power flow -- a modeling choice
made once, upstream, not a source of additional error in the
\emph{uncertainty-propagation} step) makes the affine-transformation
theorem apply without qualification.
\end{proof}

Because DC power flow is linear, $\text{flow}=
\text{PTDF}\cdot\text{injections}$, Gaussian injection uncertainty
propagates \emph{exactly} via
\begin{equation}
\Sigma_{\text{flow}} = \text{PTDF}\,\Sigma_{\text{inj}}\,\text{PTDF}^\top,
\label{eq:ppf}
\end{equation}
with overload probability $\Pr(\left|\text{flow}\right|>\text{rate}A)$ from the
two-sided Gaussian tail -- computed via the identical quantile machinery
as Theorem~\ref{thm:chance-tight}, now applied to a flow margin instead
of a reserve margin. Verified on a 3-bus test system
($\text{rate}A=[30,100,100]$, $\text{inj\_mean}=[0,-60,40]$,
$\text{inj\_cov}=\text{diag}([0,25,225])$): analytical branch-1 flow
mean/std of $26.667/6.009$\,MW matches 20{,}000-sample Monte Carlo
($26.658/6.005$) closely, with overload probability $0.2895$ (analytical)
vs.\ $0.2879$ (Monte Carlo) -- and collapses exactly to a deterministic
DC power-flow solve at zero variance. MATLAB reproduces every number to
4 decimal places.

\section{Dynamic line rating (IEEE 738)}

\begin{figure}[H]
\centering
\includegraphics[width=0.46\textwidth]{Images/flowchart_dynamic_line_rating.png}
\caption{Process flow of \pkg{dynamic\_line\_rating.py}.}
\end{figure}

\pkg{dynamic\_\hspace{0pt}line\_\hspace{0pt}rating.py}/\pkg{psat\_\hspace{0pt}dynamic\_\hspace{0pt}line\_\hspace{0pt}rating.m}
compute weather-dependent conductor ampacity from the IEEE 738
steady-state heat balance,
\begin{equation}
I_{\max} = \sqrt{\frac{q_c+q_r-q_s}{R(T_c)}},
\label{eq:dlr}
\end{equation}
with radiative loss, solar gain, and the larger of two McAdams-type
forced-convection correlations (each the IEEE 738 standard's own
published equation), while the underlying air-property coefficients
(density, viscosity, thermal conductivity as functions of film
temperature) use general ideal-gas/Sutherland's-law physics where this
project could not independently verify the standard's own specific
tables -- stated honestly rather than claimed as a verified reproduction.
Verified across three weather scenarios (static-like: 793.8\,A; cold and
windy: 2284.1\,A; hot and calm: 165.5\,A), every result's heat-balance
residual near machine precision ($\sim10^{-13}$), with the wind-angle
factor
\begin{equation}
K(\phi) = 1.194-\cos\phi+0.194\cos2\phi+0.368\sin2\phi
\label{eq:windangle}
\end{equation}
confirmed at its published boundary values ($1.0$ at $90^\circ$, $0.388$
at $0^\circ$). MATLAB reproduces every ampacity to the decimal.

\begin{proposition}[Existence and uniqueness of the IEEE 738 ampacity solution]
\label{prop:dlr-unique}
For fixed weather conditions and conductor temperature $T_c$, the
ampacity $I_{\max}$ solving $I^2R(T_c)+q_s=q_c(I)+q_r$ (resistive
heating plus any temperature-dependence of $R$ held fixed, balanced
against convective and radiative loss, net of solar gain) has a unique
non-negative solution whenever net cooling at zero current is
non-negative, because $I\mapsto I^2R(T_c)$ is strictly increasing and
convex for $I\ge0$ while the right-hand side is independent of $I$ in
this module's own formulation (convective loss $q_c$ here depends on
wind/ambient conditions and $T_c$, not on the current $I$ itself, since
$T_c$ is the \emph{fixed} maximum allowable operating temperature, not a
free variable solved jointly with $I$).
\end{proposition}
\begin{proof}
Define $g(I)=I^2R(T_c)-(q_c+q_r-q_s)$. Since $R(T_c)>0$ is fixed, $g$ is
a strictly increasing, continuous function of $I\ge0$ with
$g(0)=-(q_c+q_r-q_s)$. If $q_c+q_r-q_s\ge0$ (net cooling available),
$g(0)\le0$ and $g(I)\to\infty$ as $I\to\infty$, so by the intermediate
value theorem combined with strict monotonicity, $g$ has exactly one
root $I_{\max}\ge0$ -- which is precisely
$I_{\max}=\sqrt{(q_c+q_r-q_s)/R(T_c)}$, the closed form
the IEEE 738 heat-balance equation already uses directly rather than via iterative root-
finding. If instead $q_c+q_r-q_s<0$ (solar heating alone already
exceeds available cooling), no non-negative $I$ solves $g(I)=0$, and the
module correctly returns $I_{\max}=0$ by convention (Chapter~\ref{ch:distribution}'s
verified ``extreme heat and calm'' test case) rather than a spurious
negative or complex root, matching exactly this proposition's boundary
case.
\end{proof}

%=============================================================
\chapter{Interoperability with MATPOWER}
\label{ch:interop}
%=============================================================

\begin{figure}[H]
\centering
\includegraphics[width=0.46\textwidth]{Images/flowchart_matpower_m_parser.png}
\hfill
\includegraphics[width=0.46\textwidth]{Images/flowchart_matpower_to_psat.png}
\caption{Process flow of \pkg{matpower\_m\_parser.py} (left) and \pkg{matpower\_to\_psat.py} (right).}
\end{figure}

\pkg{matpower\_\hspace{0pt}m\_\hspace{0pt}parser.py} parses a raw MATPOWER case \texttt{.m} file
(\texttt{mpc.bus}, \texttt{mpc.gen}, \texttt{mpc.branch}, optionally
\texttt{mpc.gencost}) into a PYPOWER-compatible dictionary, closing the
gap that PYPOWER's own \texttt{loadcase()} accepts only \texttt{.mat}/
\texttt{.py} case files, not the plain-text format MATPOWER itself (and
this project's larger downloaded PEGASE/RTE cases) actually ships in.
\pkg{matpower\_\hspace{0pt}to\_\hspace{0pt}psat.py} converts a MATPOWER/PYPOWER case dictionary
into a PSAT-native data file (\texttt{Bus.con}, \texttt{Line.con},
\texttt{SW.con}, \texttt{PV.con}, \texttt{PQ.con}, optionally
\texttt{Supply.con}/\texttt{Demand.con}), using the source-verified
\pkg{con\_\hspace{0pt}schema.json} column mapping, so a 118-bus or 300-bus standard
test system can be run directly in PSAT, not only in PYPOWER. Generator
dynamic data (inertia, reactances, AVR/governor parameters) is
deliberately \emph{not} fabricated, since a MATPOWER case carries none --
converted cases are power-flow/OPF-ready but explicitly not claimed
ready for time-domain or small-signal analysis. Both utilities are
Python-side file-format infrastructure by design and have no MATLAB
``port'': the second tool's whole purpose is to write a file MATLAB
already loads directly via \texttt{runpsat}.

%=============================================================
\chapter{Verification Methodology and Consolidated Defect Log}
\label{ch:verification}
%=============================================================

\section{Cross-language parity}

Every Python module with a natural PSAT-side analogue has a MATLAB port,
each tested live against MATLAB R2026a. (MATLAB in this development
environment must be invoked with its process sandbox disabled; under the
default sandbox, \texttt{matlab -batch} fails with a license-service
communication error even though the license is valid and network
connectivity is fine -- a tool-sandbox artifact, not a real licensing
failure.) Two Python modules are intentionally \emph{not} ported:
reinforcement learning (\pkg{rl\_\hspace{0pt}env.py}) is fundamentally a
Python/Gymnasium-ecosystem capability, and document generation
(\pkg{lab\_\hspace{0pt}report\_\hspace{0pt}generator.py}) is already well served by
\pkg{python-docx}, with a parallel MATLAB implementation judged lower
value than the other extensions for the time available. The two
MATPOWER-format utilities are Python-only infrastructure by design
(\S\ref{ch:interop}).

\begin{longtable}{@{}p{5cm}p{5.5cm}p{4cm}@{}}
\caption{Python/MATLAB extension parity map.}
\label{tab:parity}\\
\toprule
Python module & MATLAB port(s) & Verification \\
\midrule
\endfirsthead
\toprule
Python module & MATLAB port(s) & Verification \\
\midrule
\endhead
\pkg{psat\_\hspace{0pt}io.py} & (reader only; no port needed) & Tested, 3 real files \\
\pkg{gas\_\hspace{0pt}network.py} & \pkg{psat\_\hspace{0pt}gas\_\hspace{0pt}network.m} & Live, matches to 2 d.p. \\
\pkg{decentralized\_\hspace{0pt}opt.py} & \pkg{psat\_\hspace{0pt}consensus\_\hspace{0pt}dispatch.m}, \pkg{psat\_\hspace{0pt}decentralized\_\hspace{0pt}opt.m} & Live, exact match \\
\pkg{stochastic\_\hspace{0pt}opf.py} & \pkg{psat\_\hspace{0pt}chance\_\hspace{0pt}constrained\_\hspace{0pt}reserve.m} & Live, exact match \\
\pkg{market\_\hspace{0pt}planning.py} & \pkg{psat\_\hspace{0pt}merit\_\hspace{0pt}order.m}, \pkg{psat\_\hspace{0pt}market\_\hspace{0pt}planning.m} & Live, exact match (after dual-sign fix) \\
\pkg{fdia\_\hspace{0pt}state\_\hspace{0pt}estimation.py} & \pkg{psat\_\hspace{0pt}fdia\_\hspace{0pt}demo.m} & Live, $J$ unchanged under attack \\
\pkg{mpc\_\hspace{0pt}export.py} & \pkg{psat\_\hspace{0pt}discretize.m}, \pkg{psat\_\hspace{0pt}mpc\_\hspace{0pt}design.m} & Live, exact LQR-gain match \\
\pkg{rl\_\hspace{0pt}env.py} & \emph{not ported (by design)} & -- \\
\pkg{visualization\hspace{0pt}.py} & \pkg{psat\_\hspace{0pt}visualize.m} & Live \\
\pkg{three\_\hspace{0pt}phase\_\hspace{0pt}distribution.py} & \pkg{psat\_\hspace{0pt}three\_\hspace{0pt}phase\_\hspace{0pt}distribution.m} & Live, 4 d.p.\ match after transpose fix \\
\pkg{short\_\hspace{0pt}circuit.py} & \pkg{psat\_\hspace{0pt}short\_\hspace{0pt}circuit.m} & Live, exact digit match \\
\pkg{battery\_\hspace{0pt}storage.py} & \pkg{psat\_\hspace{0pt}battery\_\hspace{0pt}storage.m} & Live, exact match, no bugs \\
\pkg{robust\_\hspace{0pt}opf.py} & \pkg{psat\_\hspace{0pt}robust\_\hspace{0pt}opf.m}, \pkg{psat\_\hspace{0pt}robust\_\hspace{0pt}multi\_\hspace{0pt}uncertainty.m} & Live, exact match \\
\pkg{cascading\_\hspace{0pt}failure.py} & \pkg{psat\_\hspace{0pt}cascading\_\hspace{0pt}failure.m} & Live, exact match \\
\pkg{wide\_\hspace{0pt}area\_\hspace{0pt}control.py} & \pkg{psat\_\hspace{0pt}wide\_\hspace{0pt}area\_\hspace{0pt}control.m} & Live, both instability and fix reproduced \\
\pkg{attack\_\hspace{0pt}surface.py} & \pkg{psat\_\hspace{0pt}attack\_\hspace{0pt}surface.m} & Live on WSCC-9 \\
\pkg{gnn\_\hspace{0pt}power\_\hspace{0pt}flow.py} & \pkg{psat\_\hspace{0pt}gnn\_\hspace{0pt}power\_\hspace{0pt}flow.m} & Live, trained on 200 real PSAT NR solves \\
\pkg{most\_\hspace{0pt}scheduling.py} & \pkg{psat\_\hspace{0pt}most.m} & Live, cost match to the penny \\
\pkg{stability\_\hspace{0pt}constrained\_\hspace{0pt}opf.py} & \pkg{psat\_\hspace{0pt}stability\_\hspace{0pt}constrained\_\hspace{0pt}opf.m} & Live, exact match incl.\ 90$^\circ$ threshold \\
\pkg{probabilistic\_\hspace{0pt}power\_\hspace{0pt}flow.py} & \pkg{psat\_\hspace{0pt}probabilistic\_\hspace{0pt}power\_\hspace{0pt}flow.m} & Live, 4 d.p.\ match \\
\pkg{dynamic\_\hspace{0pt}line\_\hspace{0pt}rating.py} & \pkg{psat\_\hspace{0pt}dynamic\_\hspace{0pt}line\_\hspace{0pt}rating.m} & Live, exact digit match \\
\pkg{lab\_\hspace{0pt}report\_\hspace{0pt}generator.py} & \emph{not ported (by design)} & -- \\
\pkg{matpower\_\hspace{0pt}m\_\hspace{0pt}parser.py} & \emph{(Python-only infrastructure)} & -- \\
\pkg{matpower\_\hspace{0pt}to\_\hspace{0pt}psat.py} & \emph{(Python-only infrastructure)} & -- \\
\emph{(none -- new device)} & \pkg{@ESSclass} & Live, PSAT core integration \\
\emph{(none -- new device)} & \pkg{@GFMclass} & Live, PSAT core integration \\
\bottomrule
\end{longtable}

\section{Consolidated defect log}

Table~\ref{tab:bugs} lists every genuine defect found and fixed during
development, in rough chronological/thematic order, kept deliberately
blunt about root cause since the recurring \emph{categories} are the
real lesson: sign-convention errors (3 independent occurrences across
two languages and two solvers), the slack-bus-left-unconstrained DC-OPF
formulation error (documented once, then independently re-triggered in
two later, unrelated modules), and MATLAB-specific pitfalls
(\texttt{'} as conjugate transpose, nested functions closing over outer
data, local functions not externally callable) that would trip up any
future port in this ecosystem.

\begin{longtable}{@{}p{3.6cm}p{11.9cm}@{}}
\caption{Consolidated defect log.}
\label{tab:bugs}\\
\toprule
Module & Defect and fix \\
\midrule
\endfirsthead
\toprule
Module & Defect and fix \\
\midrule
\endhead
\pkg{psat\_\hspace{0pt}results.m} & \texttt{DAE.y} orders angles before voltage magnitude, opposite of a first assumption; caught via a known slack-voltage value. \\
\pkg{psat\_\hspace{0pt}n1\_\hspace{0pt}classify.m}, \pkg{psat\_\hspace{0pt}secure\_\hspace{0pt}atc.m} & Editing \texttt{Line.con} in memory and re-calling \texttt{runpsat('pf')} has no effect (PSAT re-reads from disk every call); fixed by writing a physically modified file to disk. \\
\pkg{psat\_\hspace{0pt}write\_\hspace{0pt}line\_\hspace{0pt}outage\_\hspace{0pt}file.m} & Three sub-bugs: overwrote the wrong trailing column on a 15-field row; a MATLAB matrix literal required uniform row width; the correct default pad value for an omitted status column is 1 (in service), not 0. \\
\pkg{psat\_\hspace{0pt}export\_\hspace{0pt}ss.m} & No PSAT routine \texttt{fex\_abcd} exists; the real routine is \texttt{fm\_abcd}, with a different (global-struct) calling convention entirely. \\
\pkg{psat\_\hspace{0pt}ss\_\hspace{0pt}loader.py} & MATLAB's \texttt{A(:)} flattens column-major; a first draft reshaped row-major, silently transposing non-square matrices. \\
\pkg{@ESSclass}, \pkg{@GFMclass} & A sign-convention error in how generation is injected into PSAT's power-balance equation caused a genuine small-signal instability (eigenvalue $+25.9$), found via \texttt{sssa}, and found to also affect the already-shipped \pkg{@ESSclass}. \\
\pkg{decentralized\_\hspace{0pt}opt.py} & Fixed step size $0.5$ caused $\lambda$ to diverge; a coarse grid search (not Brent's method) in the ADMM inner problem left a quantization bias; stopping on aggregate balance alone can mask non-consensus among agents. \\
\pkg{market\_\hspace{0pt}planning.py} (CVXPY) & Raw CVXPY dual for the balance equality returned the \emph{negative} of the correct LMP. \\
\pkg{psat\_\hspace{0pt}market\_\hspace{0pt}planning.m} (\texttt{linprog}) & The identical dual-sign bug, independently re-discovered in a different language and solver; fixed the same way (negate, then cross-check against the known-correct uncongested price). \\
\pkg{most\_\hspace{0pt}scheduling.py} & A quadratic objective is silently rejected by both MILP solvers used; a genuine LODF-formula bug (spurious reactance factor) let a ``security-constrained'' dispatch violate its own limit by 18\,MW. \\
\pkg{most\_\hspace{0pt}scheduling.py}'s own test script & An independently-written cvxpy cross-check reproduced the slack-bus-left-unconstrained DC-OPF bug already documented from this course's own Week~6 development -- found again, in a brand-new script, months later. \\
\pkg{stability\_\hspace{0pt}constrained\_\hspace{0pt}opf.py} & \texttt{economic\_dispatch()} crashed with \texttt{TypeError} on genuine infeasibility instead of reporting it; fixed with an explicit \texttt{None} check. \\
\pkg{psat\_\hspace{0pt}stability\_\hspace{0pt}constrained\_\hspace{0pt}opf.m} & A correct but cosmetically noisy \texttt{quadprog} call on an all-zero Hessian (a pure LP); switched to \texttt{linprog}, identical results, no warning. \\
\pkg{psat\_\hspace{0pt}gas\_\hspace{0pt}network.m} & MATLAB nested functions closing over outer data produced silent \texttt{NaN}/singular-Jacobian results despite identical arithmetic working as a plain script; fixed with ordinary local functions and explicit arguments. \\
\pkg{psat\_\hspace{0pt}three\_\hspace{0pt}phase\_\hspace{0pt}distribution.m} & MATLAB's \texttt{'} is conjugate transpose for complex arrays, silently swapping phases B and C; produced a fully converged, wrong answer with no error. \\
\pkg{test\_\hspace{0pt}psat\_\hspace{0pt}dlr.m} & First draft called a MATLAB \emph{local} function from outside its own file (not externally callable); fixed by testing the externally observable behaviour instead. \\
\pkg{gnn\_\hspace{0pt}power\_\hspace{0pt}flow.py} & An un-normalized shared loss across two outputs of very different magnitude let one output dominate every gradient step, producing a model worse than a trivial baseline on the other output. \\
\pkg{wide\_\hspace{0pt}area\_\hspace{0pt}control.py} & Applying the consensus weight matrix once per physical step, rather than repeatedly, is unstable (diverges to tens of thousands of Hz), not merely slow. \\
\bottomrule
\end{longtable}

%=============================================================
\chapter{Known Limitations}
\label{ch:limitations}
%=============================================================

\begin{itemize}
\item \textbf{MATLAB sandbox dependency.} MATLAB/PSAT must be invoked
  with the development environment's process sandbox disabled; this is a
  tooling artifact of this specific environment, not a PSAT or MATLAB
  licensing limitation.
\item \textbf{Device classes, v1 scope.} \pkg{@ESSclass} and
  \pkg{@GFMclass} both omit reactive-power support and use a fixed
  (non-time-varying, non-droop) reference signal; natural v2 extensions
  are noted in each class's own README.
\item \textbf{\pkg{dynamic\_\hspace{0pt}line\_\hspace{0pt}rating.py}}'s air-property coefficients
  are general ideal-gas physics, not independently verified against IEEE
  738's own specific published tables; results should be treated as
  representative, not as a drop-in certified calculation.
\item \textbf{\pkg{gnn\_\hspace{0pt}power\_\hspace{0pt}flow.py}} does not generalize across
  network topologies; a surrogate trained on one topology must be
  retrained for another.
\item \textbf{\pkg{most\_\hspace{0pt}scheduling.py}} and its MATLAB port are
  DC-network-only (matching MOST's own current released scope); ramping
  \emph{costs} (as distinct from ramping \emph{limits}) are not
  separately priced.
\item \textbf{\pkg{matpower\_\hspace{0pt}to\_\hspace{0pt}psat.py}} converts power-flow/OPF data
  only; it does not fabricate dynamic (inertia, AVR, governor) data, so
  converted cases are not ready for time-domain or small-signal analysis
  without further, separately sourced data.
\item \textbf{No live-Canvas or production-grid validation.} Every
  numeric result in this report was verified against an independent
  method, solver, or closed-form result within this development
  environment; none of these tools have been validated against a live
  utility SCADA/EMS system or a production grid deployment.
\end{itemize}

%=============================================================
\chapter{Conclusion and Future Work}
\label{ch:conclusion}
%=============================================================

This report has documented 25 Python modules, 31 MATLAB functions, and
two genuine new PSAT device classes, closing every gap identified in
this project's own \emph{MATPOWER vs.\ PSAT Capability Comparison
Report} and in two subsequent rounds of ``what else would be valuable
for modern power systems research'' review. Every capability was
verified against an independent method before acceptance, and every
genuine defect found along the way -- 19 of them, spanning sign
conventions, MATLAB-specific semantics, and optimization-formulation
pitfalls -- is recorded in Table~\ref{tab:bugs} rather than silently
corrected and forgotten, because several of these defect \emph{classes}
recurred more than once and are worth a future developer's attention
before, not after, they recur a third time.

Natural next steps, in rough priority order: (1) time-varying/droop
reference signals and reactive-power support for \pkg{@ESSclass}/
\pkg{@GFMclass}; (2) an AC (rather than DC-only) extension to
\pkg{most\_\hspace{0pt}scheduling.py}, mirroring MATPOWER MOST's own in-development
roadmap; (3) a verified reproduction of IEEE 738's own specific
air-property tables for \pkg{dynamic\_\hspace{0pt}line\_\hspace{0pt}rating.py}; (4) a
cross-topology-generalizing GNN architecture (e.g.\ explicit
per-unit/topology-invariant feature engineering) for
\pkg{gnn\_\hspace{0pt}power\_\hspace{0pt}flow.py}; and (5) validation of the Canvas QTI and
Mentimeter deliverables documented elsewhere in this project against a
live Canvas import, the one verification step this development
environment cannot perform on its own.

\begin{thebibliography}{9}
\bibitem{milano} F.~Milano, \emph{Power System Modelling and Scripting},
  Springer, 2010. (PSAT 2.1.11 reference manual and source.)
\bibitem{zimmerman} R.D.~Zimmerman, C.E.~Murillo-S\'anchez, R.J.~Thomas,
  ``MATPOWER: Steady-State Operations, Planning, and Analysis Tools for
  Power Systems Research and Education,'' \emph{IEEE Trans.\ Power
  Syst.}, 2011.
\bibitem{most} R.D.~Zimmerman, C.E.~Murillo-S\'anchez, ``MATPOWER
  Optimal Scheduling Tool (MOST) User's Manual.''
\bibitem{bertsimas} D.~Bertsimas, M.~Sim, ``The Price of Robustness,''
  \emph{Operations Research}, 52(1), 2004.
\bibitem{boyd} S.~Boyd, N.~Parikh, E.~Chu, B.~Peleato, J.~Eckstein,
  ``Distributed Optimization and Statistical Learning via the Alternating
  Direction Method of Multipliers,'' \emph{Foundations and Trends in
  Machine Learning}, 2011.
\bibitem{liu} Y.~Liu, P.~Ning, M.K.~Reiter, ``False Data Injection
  Attacks against State Estimation in Electric Power Grids,'' \emph{ACM
  CCS}, 2009.
\bibitem{ieee738} IEEE Standard 738-2012, \emph{IEEE Standard for
  Calculating the Current-Temperature Relationship of Bare Overhead
  Conductors}.
\bibitem{glover} J.D.~Glover, T.J.~Overbye, M.S.~Sarma, \emph{Power
  System Analysis and Design}, Cengage Learning, 6th ed., 2016.
\bibitem{kundur} P.~Kundur, \emph{Power System Stability and Control},
  McGraw-Hill, 1994.
\bibitem{boydconvex} S.~Boyd, L.~Vandenberghe, \emph{Convex
  Optimization}, Cambridge University Press, 2004.
\bibitem{bertsekasDP} D.P.~Bertsekas, \emph{Dynamic Programming and
  Optimal Control}, Vol.\ I--II, Athena Scientific, 4th ed., 2017.
\bibitem{bertsekasnlp} D.P.~Bertsekas, \emph{Nonlinear Programming},
  Athena Scientific, 3rd ed., 2016.
\bibitem{bertsimastsitsiklis} D.~Bertsimas, J.N.~Tsitsiklis,
  \emph{Introduction to Linear Optimization}, Athena Scientific, 1997.
\bibitem{hornjohnson} R.A.~Horn, C.R.~Johnson, \emph{Matrix Analysis},
  Cambridge University Press, 2nd ed., 2012.
\bibitem{golubvanloan} G.H.~Golub, C.F.~Van Loan, \emph{Matrix
  Computations}, Johns Hopkins University Press, 4th ed., 2013.
\bibitem{xiaoboydconsensus} L.~Xiao, S.~Boyd, ``Fast Linear Iterations
  for Distributed Averaging,'' \emph{Systems \& Control Letters}, 53(1),
  2004.
\bibitem{casellaberger} G.~Casella, R.L.~Berger, \emph{Statistical
  Inference}, Duxbury/Cengage Learning, 2nd ed., 2002.
\bibitem{aburexposito} A.~Abur, A.G.~Exp\'osito, \emph{Power System
  State Estimation: Theory and Implementation}, CRC Press, 2004.
\bibitem{andersonmoore} B.D.O.~Anderson, J.B.~Moore, \emph{Optimal
  Control: Linear Quadratic Methods}, Prentice Hall, 1990 (Dover
  reprint, 2007).
\bibitem{nocedalwright} J.~Nocedal, S.J.~Wright, \emph{Numerical
  Optimization}, Springer, 2nd ed., 2006.
\bibitem{birgelouveaux} J.R.~Birge, F.~Louveaux, \emph{Introduction to
  Stochastic Programming}, Springer, 2nd ed., 2011.
\bibitem{kersting} W.H.~Kersting, \emph{Distribution System Modeling
  and Analysis}, CRC Press, 4th ed., 2017.
\bibitem{chengshirmohammadi} C.S.~Cheng, D.~Shirmohammadi, ``A
  Three-Phase Power Flow Method for Real-Time Distribution System
  Analysis,'' \emph{IEEE Trans.\ Power Syst.}, 10(2), 1995.
\bibitem{graingerstevenson} J.J.~Grainger, W.D.~Stevenson Jr.,
  \emph{Power System Analysis}, McGraw-Hill, 1994.
\bibitem{woodwollenberg} A.J.~Wood, B.F.~Wollenberg, G.B.~Sheble,
  \emph{Power Generation, Operation, and Control}, Wiley, 3rd ed., 2013.
\bibitem{machowski} J.~Machowski, J.W.~Bialek, J.R.~Bumby, \emph{Power
  System Dynamics: Stability and Control}, Wiley, 2nd ed., 2008.
\end{thebibliography}

\end{document}
